% ==============================================================================
% Manuscript: Hybrid Deep Learning Ensemble for Knee Osteoarthritis Severity
% Grading from Radiographic Images with Ordinal-Aware Optimization and
% Multi-Method Explainability
% ==============================================================================
\documentclass[10pt,journal]{IEEEtran}

\usepackage{amsmath}
\usepackage{amssymb}
\usepackage{graphicx}
\graphicspath{{./}{./figures/}{./images/}{./dl-images/}}
\DeclareGraphicsExtensions{.jpeg,.jpg,.png,.pdf}
\usepackage{booktabs}
\usepackage{multirow}
\usepackage{array}
\usepackage{url}
\usepackage{xcolor}
\usepackage{cite}
\usepackage{algorithm}
\usepackage{algorithmic}
\usepackage{bm}
\usepackage{hyperref}
\usepackage{tabularx}
\usepackage{enumitem}
\usepackage{float}
\usepackage{dblfloatfix}
\usepackage[section]{placeins}

% Double-column float placement optimization to keep figures inline with text
\setcounter{topnumber}{4}
\setcounter{bottomnumber}{4}
\setcounter{totalnumber}{10}
\setcounter{dbltopnumber}{4}
\renewcommand{\topfraction}{0.99}
\renewcommand{\bottomfraction}{0.99}
\renewcommand{\textfraction}{0.01}
\renewcommand{\floatpagefraction}{0.50}
\renewcommand{\dbltopfraction}{0.99}
\renewcommand{\dblfloatpagefraction}{0.50}

\hypersetup{
    colorlinks=true,
    linkcolor=blue,
    citecolor=blue,
    urlcolor=blue
}

\newcommand{\TODO}[1]{\textcolor{red}{\textbf{[TO BE REPLACED: #1]}}}
\newcommand{\NOTPROVIDED}[1]{\textcolor{orange}{\textbf{[EXPERIMENTAL DETAIL NOT PROVIDED: #1]}}}

\begin{document}

\title{Hybrid Deep Learning Ensemble with Ordinal-Aware Loss for Automated Knee Osteoarthritis Severity Grading from Radiographic Images}

\author{
\begin{minipage}[t]{0.31\textwidth}
\centering
\textbf{Setty Bhavithav}$^{1}$\\
$^{}$Department of Computer Science and Engineering\\
Amrita Vishwa Vidyapeetham, Amaravati, Andhra Pradesh, India\\
E-mail: \href{mailto:settybhavithav@gmail.com}{settybhavithav@gmail.com}\\
ORCID: \href{https://orcid.org/0009-0009-1065-6154}{0009-0009-1065-6154}
\end{minipage}
\hfill
\begin{minipage}[t]{0.31\textwidth}
\centering
\textbf{MNGR Bharadwaj}$^{2}$\\
$^{}$Department of Computer Science and Engineering\\
Amrita Vishwa Vidyapeetham, Amaravati, Andhra Pradesh, India\\
E-mail: \href{mailto:bharadwaj4806@gmail.com}{bharadwaj4806@gmail.com}\\
ORCID: \href{https://orcid.org/0009-0000-1065-8701}{0009-0000-1065-8701}
\end{minipage}
\hfill
\begin{minipage}[t]{0.31\textwidth}
\centering
\textbf{Kelam Bhargav}$^{3}$\\
$^{}$Department of Computer Science and Engineering\\
Amrita Vishwa Vidyapeetham, Amaravati, Andhra Pradesh, India\\
E-mail: \href{mailto:barghavkelam04@gmail.com}{barghavkelam04@gmail.com}\\
ORCID: \href{https://orcid.org/0009-0001-7333-8104}{0009-0001-7333-8104}
\end{minipage}
}

\maketitle

% ==============================================================================
% ABSTRACT
% ==============================================================================
\begin{abstract}
Radiographic grading of knee osteoarthritis (KOA) under the semi-quantitative Kellgren--Lawrence (KL) scale is hindered by high inter-reader variability, potential data leakage across bilateral examinations, and the failure of nominal loss objectives to respect ordinal disease stages. This study presents a triage-oriented hybrid deep learning architecture that integrates three complementary visual paradigms: ConvNeXt-Base for subchondral bone trabecular texture, Swin Transformer for shifted-window joint compartment context, and self-supervised DINO Vision Transformer for invariant structural boundary priors. The five nominal KL grades are operationally grouped into three study-defined triage tiers: Low/No OA (KL~0--1), Moderate OA (KL~2--3), and Severe OA (KL~4). To mitigate non-adjacent diagnostic errors, a distance-aware ordinal penalty loss is integrated into optimization to penalize misclassifications proportionally to continuous metric distance in label space. Evaluated across a core cohort of 8,260 radiographs from 4,130 unique patients with verified patient-level isolation ($|\mathcal{P}_{\text{Train}} \cap \mathcal{P}_{\text{Test}}| = 0$), the primary held-out test partition ($N = 1{,}656$ radiographs, 828 subjects) achieved a static single-pass accuracy of 83.44\%, Quadratic Weighted Kappa (QWK) of 0.721, and macro ROC-AUC of 0.943, advancing to 85.08\% accuracy, 0.748 QWK, and 0.954 macro ROC-AUC under horizontal-flip test-time augmentation (TTA), with 100\% empirical sensitivity for Severe OA (51/51) and zero non-adjacent omissions. On an expanded 4,008-image multi-partition stress-testing cohort, the Hybrid model attained 90.92\% accuracy, a QWK of 0.853, and a macro ROC-AUC of 0.980, with zero observed non-adjacent omissions across all 122 Severe cases. Interpretability audits across gradient and perturbation methods show consistent qualitative localization along the tibiofemoral joint margins, supporting the potential use of the framework as an assistive radiographic triage tool.
\end{abstract}

\begin{IEEEkeywords}
Knee osteoarthritis, deep learning, vision transformer, ordinal classification, explainable artificial intelligence, computer-aided diagnosis.
\end{IEEEkeywords}

% ==============================================================================
% 1. INTRODUCTION
% ==============================================================================
\section{Introduction}
\label{sec:introduction}

Knee osteoarthritis (KOA) is the primary degenerative joint disorder worldwide, affecting over 300 million individuals and imposing severe physical disability and socioeconomic burden on healthcare systems \cite{hunter2019oa, loeser2012oa}. The pathology involves progressive degradation of articular cartilage, osteophytic proliferation along bone margins, subchondral sclerosis, and progressive loss of joint-space width \cite{cross2014global, woolf2003burden}. Plain weight-bearing fixed-flexion radiography represents the standard frontline imaging modality for semi-quantitative assessment of KOA owing to its cost efficiency, immediate accessibility, and minimal radiation exposure \cite{braun2012diagnosis}. Semiquantitative staging has relied on the Kellgren--Lawrence (KL) classification since 1957, scoring radiographs from Grade~0 (normal) to Grade~4 (severe) \cite{kellgren1957radiological}.

Despite widespread clinical adoption, manual KL interpretation suffers from substantial subjectivity. Inter-reader agreement among board-certified musculoskeletal radiologists frequently yields weighted kappa ($\kappa$) coefficients between 0.50 and 0.72, with marked divergence along transitional boundaries \cite{wright2014interobserver, sheehy2015validity}. Specifically, differentiating doubtful joint narrowing (KL~1) from definitive early osteophytic lipping (KL~2) represents a common diagnostic dilemma that impacts timely intervention \cite{culvenor2019oa}. Automated computer-aided diagnostic algorithms powered by deep neural networks hold promise for assisting severity grading and standardizing diagnostic workflows \cite{tiulpin2018automatic, thomas2020automated, chen2019fully}.

Nevertheless, an appraisal of the literature highlights three methodological vulnerabilities. First, several investigations partition benchmark datasets at the image level rather than the patient level. Because bilateral radiographs from the same individual share genetic bone morphology, anatomical habitus, and exposure parameters, random image splitting can permit correlated observations from the same patient to appear across partitions, potentially leading to optimistic estimates of generalization performance. Second, prevailing models optimize standard nominal cross-entropy loss, which penalizes all classification errors uniformly. Clinically, confusing Severe OA (KL~4, advanced joint failure) with Low/No OA (KL~0--1) represents a hazardous omission, whereas misclassifying an adjacent transition (KL~1 vs.\ KL~2) carries manageable clinical consequences \cite{antony2016quantifying}. Third, single-backbone architectures exhibit intrinsic inductive bias trade-offs: standard convolutional neural networks (CNNs) prioritize localized high-frequency edge textures at the expense of long-range spatial context, whereas vision transformers generally rely more heavily on large-scale pretraining to learn useful visual representations when applied to moderately sized medical datasets.

To address these challenges, this study introduces an automated, leakage-controlled diagnostic framework combining convolutional, shifted-window self-attention, and self-supervised visual priors. The primary contributions of this work are three-fold:

\begin{enumerate}[leftmargin=*, label=\arabic*.]
\item \textbf{Tri-Paradigm Representation Synthesis:} We formulate a multi-branch late-fusion architecture that concurrently extracts three complementary visual representations: ConvNeXt-Base for localized subchondral bone texture, Swin Transformer for hierarchical shifted-window joint compartment context, and DINO ViT for self-supervised structural boundary priors. This unified feature space mitigates the receptive field limitations of single-backbone KOA classifiers.

\item \textbf{Distance-Aware Ordinal-Penalized Optimization:} We develop a composite loss function combining label-smoothed categorical cross-entropy with a continuous expected-value metric penalty in label space. By explicitly penalizing the magnitude of ordinal grade displacement, the optimization objective suppresses non-adjacent prediction errors, mitigating severe misclassifications while preserving sensitivity along borderline intermediate stages.

\item \textbf{Leakage-Controlled Evaluation and Multi-Branch Auditing:} Across 8,260 core radiographs from 4,130 unique subjects, we enforce verified patient-level isolation across training (2,889 patients, 5,778 images), validation (413 patients, 826 images), and primary held-out test (828 patients, 1,656 images) partitions. On the untouched primary test set ($N = 1{,}656$), the Hybrid architecture attains 83.44\% static accuracy (0.721 QWK, 0.943 AUC), advancing to 85.08\% accuracy (0.748 QWK, 0.954 AUC) with horizontal-flip test-time augmentation (TTA). On an expanded 4,008-image multi-partition stress-testing cohort, the framework reaches 90.92\% accuracy, 0.853 QWK, and 0.980 macro ROC-AUC, achieving 100\% sensitivity on advanced disease ($N = 122$) with zero observed extreme reversals into Low/No OA. Multi-method explainability audits across gradient and perturbation methods show attribution patterns consistent with relevant anatomical regions along the tibiofemoral joint margins.
\end{enumerate}

The remainder of this article is structured as follows. Section~\ref{sec:literature} surveys related literature and delineates the research gap. Section~\ref{sec:methodology} describes the materials, mathematical formulation, and explainability framework. Section~\ref{sec:results} presents quantitative benchmark results, ablation dynamics, and visual attributions. Section~\ref{sec:discussion} contextualizes clinical decision support, examines limitations, and outlines future directions. Section~\ref{sec:conclusion} concludes the paper.


% ==============================================================================
% 2. LITERATURE REVIEW
% ==============================================================================
\section{Literature Review}
\label{sec:literature}

\subsection{Clinical Pathophysiology and Conventional Radiographic Grading}
\label{subsec:clinical_background}

Knee osteoarthritis is a whole-organ joint disorder involving dynamic interactions between mechanical stress, subchondral bone turnover, and localized synovial inflammation \cite{loeser2012oa, hunter2019oa}. Primary radiographic manifestations comprise marginal osteophytosis (calcified spurs protruding from articular margins), non-uniform joint-space narrowing (JSN) indicative of cartilage loss, subchondral eburnation (increased bone mineral density under pressure areas), and subchondral pseudocysts \cite{braun2012diagnosis}. The Kellgren--Lawrence scale categorizes these hallmarks across five distinct grades: Grade~0 (absence of OA features); Grade~1 (doubtful joint-space narrowing and possible osteophytic lipping); Grade~2 (definite osteophytes with unimpaired or unimpaired-to-doubtful joint space); Grade~3 (moderate osteophytes, marked JSN, some sclerosis, and possible bony contour deformity); and Grade~4 (large osteophytes, marked JSN progressing to bone-on-bone contact, severe sclerosis, and definite deformity) \cite{kellgren1957radiological}.

In clinical orthopedic workflows, these five nominal grades collapse into actionable therapeutic categories: asymptomatic or preventive care (KL~0--1), conservative pharmacotherapy, physiotherapy, and intra-articular injections (KL~2--3), and total knee arthroplasty consideration (KL~4) \cite{zhang2010eular, carr2012knee}. Manual assessment requires experienced radiologists to subjectively synthesize these multi-focal signs from two-dimensional projections, resulting in substantial inter-reader variability \cite{schiphof2008differences, wright2014interobserver}.

\subsection{Evolution of Computational KOA Grading}
\label{subsec:evolution}

Early computer-assisted approaches relied on handcrafted texture descriptors, active shape models, and fractal dimension analyses of subchondral trabecular bone \cite{shamir2009early, oka2008fully, thomson2015automated}. Although mechanistically interpretable, these feature-engineering techniques exhibited high sensitivity to exposure variation, beam angle tilt, and manual region-of-interest cropping, hindering generalized deployment across multicenter imaging hardware.

The advent of deep convolutional neural networks established end-to-end representation learning directly from raw radiograph pixels. Antony~\textit{et~al.} deployed VGG-16 networks on the Osteoarthritis Initiative (OAI) repository, achieving 57.0\% multi-class accuracy for five-class KL grading \cite{antony2016quantifying}. Subsequent refinements incorporated automated knee joint detection as a preprocessing stage, achieving incremental gains through regional localization \cite{antony2017automatic, bayramoglu2021adaptive}. Tiulpin~\textit{et~al.} introduced a bilateral Siamese network architecture capable of processing paired knee patches concurrently, reaching 66.7\% accuracy and a quadratic kappa of 0.83 on the OAI benchmark \cite{tiulpin2018automatic}. Chen~\textit{et~al.} incorporated ordinal cross-entropy formulations into dense networks, demonstrating that encoding monotonicity improved classification stability across intermediate grades \cite{chen2019fully}. Swiecicki~\textit{et~al.} demonstrated that fusing posteroanterior and lateral radiographic projections using deep convolutional neural networks yielded 71.9\% accuracy and a quadratic weighted kappa of 0.91 on patient-stratified MOST cohorts, matching board-certified musculoskeletal radiologist consensus \cite{swiecicki2021deep}.

More recently, attention-based vision architectures have been explored for musculoskeletal imaging. Swin Transformer introduced shifted-window multi-head self-attention, achieving linear computational complexity while preserving multi-scale spatial receptive fields \cite{liu2021swin}. In parallel, ConvNeXt modernized convolutional blocks via depthwise separable kernels, inverted bottlenecks, and LayerNorm operations, proving that pure CNNs can match or exceed transformer representations on visual benchmarks \cite{liu2022convnext}. In parallel developments within representation learning, self-supervised distillation via DINO demonstrates that vision transformers trained without human annotations acquire emergent attention maps aligned with intrinsic anatomical contours and articular boundaries \cite{caron2021dino}.

\subsection{Visual Foundation Models and Inductive Biases in Musculoskeletal Imaging}
\label{subsec:transformers_convnets}

The translation of deep learning to orthopedic radiograph interpretation has historically been constrained by the inductive biases of underlying neural architectures. Traditional convolutional neural networks (CNNs), such as ResNet \cite{he2016deep}, DenseNet \cite{huang2017densely}, and EfficientNet \cite{tan2019efficientnet}, rely on local receptive fields governed by translational equivariance and spatial weight sharing. While this inductive bias confers high parameter efficiency and rapid convergence on high-frequency textural patterns---such as subchondral bone sclerosis, trabecular disruption, and marginal osteophytes---it inherently restricts the network's capacity to synthesize global spatial interactions across non-contiguous anatomical structures. In knee radiography, accurate staging requires comprehensive evaluation of coronal alignment, comparing joint space height between the medial and lateral compartments simultaneously.

To transcend the localized receptive fields of CNNs, Vision Transformers (ViTs) introduce global self-attention mechanisms, modeling pairwise interactions across flattened image patches \cite{liu2021swin}. However, standard ViT architectures lack the intrinsic translational equivariance of convolutional operators and exhibit quadratic computational complexity $\mathcal{O}(N^2)$ with respect to token count $N$. Consequently, standard ViTs require massive pretraining data to avoid severe overfitting when applied to moderately sized medical datasets \cite{raghu2019transfusion, tajbakhsh2016convolutional}. To mitigate these computational bottlenecks, the Swin Transformer \cite{liu2021swin} establishes hierarchical representation learning through shifted-window multi-head self-attention (SW-MSA). By limiting self-attention to non-overlapping local windows while enabling cross-window connections via cyclic shifting, Swin achieves linear computational complexity $\mathcal{O}(N)$ relative to image dimensions while preserving multi-scale hierarchical feature maps essential for compartment-level context.

Concurrently, modern convolutional designs have evolved to bridge the performance gap with transformers. ConvNeXt \cite{liu2022convnext} modernizes the standard convolutional block by adopting depthwise separable $7 \times 7$ convolutions, inverted bottleneck channel ratios, LayerNorm instead of BatchNorm, and Gaussian Error Linear Units (GELU). These architectural refinements substantially widen effective receptive fields, enabling ConvNeXt to capture localized trabecular micro-architecture with competitive computational efficiency. In parallel, self-supervised distillation paradigms, epitomized by DINO \cite{caron2021dino}, demonstrate that vision transformers trained without human annotations via student-teacher momentum centering acquire emergent attention maps that intrinsically segment anatomical structures and boundary contours. By orchestrating a tri-paradigm synthesis uniting ConvNeXt (local trabecular texture), Swin (inter-compartmental relational context), and DINO ViT (self-supervised invariant structural geometry), multi-branch fusion circumvents the inductive limitations of single-backbone medical CAD systems.

\subsection{Distance-Aware and Ordinal Objective Formulations in Medical Imaging}
\label{subsec:ordinal_formulations}

Conventional medical image classifiers typically minimize standard categorical cross-entropy (CE) loss, which implicitly treats disease categories as nominal and mutually orthogonal classes. Under standard CE, a severe diagnostic error (e.g., misclassifying Severe OA [KL~4] as Low/No OA [KL~0--1]) incurs an identical penalty to an adjacent, clinically manageable boundary disagreement (e.g., classifying Low/No OA [KL~0--1] as Moderate OA [KL~2--3]) \cite{antony2016quantifying}. In clinical musculoskeletal practice, Kellgren--Lawrence grading establishes an ordered radiographic severity scale reflecting progressively greater structural joint degeneration, characterized by advancing articular cartilage loss, joint-space narrowing, and subchondral bone deformation \cite{kellgren1957radiological, emrani2008joint}. Misclassifying a patient across non-adjacent stages compromises patient safety, either causing catastrophic delays in surgical intervention or prompting unnecessary invasive therapies \cite{carr2012knee, zhang2010eular}.

To address ordinality, several methodological strategies have emerged in the medical imaging literature. Cost-sensitive cross-entropy introduces static asymmetric weight matrices to penalize off-diagonal errors; however, these static coefficients distort gradient magnitudes and impair early training convergence. Cumulative link and proportional odds formulations decompose the multi-class task into a series of binary threshold subproblems \cite{chen2019fully}, yet they enforce rigid parallel slope assumptions that struggle with severe class imbalance \cite{johnson2019survey}. In contrast, framing ordinal progression through expected-value rank penalties regularizes the continuous manifold directly in label space \cite{chen2019fully}. By computing the continuous mathematical expectation of predicted grades and penalizing absolute metric displacement from the true ordinal rank, the optimization landscape dynamically penalizes non-adjacent transitions while allowing smooth probabilistic transitions along ambiguous boundaries.

\subsection{Identification of the Research Gap}
\label{subsec:research_gap}

Despite extensive prior work, several foundational limitations remain unaddressed:
\begin{enumerate}[leftmargin=*, label=(\roman*)]
\item \textit{Architectural Isolation:} Many prior KOA studies rely on single-family architectures---either convolutional networks or vision transformers---without harnessing the synergistic interplay between local texture encoding, hierarchical window attention, and self-supervised structural priors.
\item \textit{Nominal Objective Functions:} Most existing KOA classification models optimize standard cross-entropy, which treats misclassifications nominally and fails to penalize the clinical hazard of extreme grade reversals (e.g., misdiagnosing KL~4 as KL~0).
\item \textit{Data Leakage Vulnerabilities:} Many published pipelines partition datasets across images rather than patient identifiers, exposing their evaluation to bilateral knee leakage and overestimating real-world clinical generalizability.
\item \textit{Opaque Decision Processes:} Comprehensive interpretability comparing gradient attributions and perturbation-based superpixel models across distinct architectural branches remains limited in the KOA literature reviewed.
\end{enumerate}


% ==============================================================================
% 3. MATERIALS AND METHODOLOGY
% ==============================================================================
\section{Materials and Methodology}
\label{sec:methodology}

\subsection{Dataset Cohort and Leakage-Free Stratification}
\label{subsec:dataset}

The radiographic cohort analyzed in this study is sourced from the Knee Osteoarthritis Severity Grading Dataset published by Chen~\textit{et~al.} on Mendeley Data (Version~1, DOI: 10.17632/56rmx5bjcr.1) \cite{chen2018dataset}, which compiles bilateral weight-bearing knee radiographs acquired under standardized protocols from the Osteoarthritis Initiative (OAI) longitudinal study \cite{nevitt2006oai}. The complete repository comprises 9,786 digitized radiographs graded according to Kellgren--Lawrence criteria.

To eliminate data leakage arising from correlated bilateral joint geometry and shared subject habitus, patient identification codes (PIDs) were systematically extracted from image filenames using regular expression parsing (\texttt{\textasciicircum([0-9]+)[LR]?\textbackslash.(?:png|jpg|jpeg)\$}). The leading numeric sequence uniquely identifies the subject, while the trailing alphabetical token denotes knee laterality (`L' for left, `R' for right). This parsing isolated exactly 4,130 unique subjects across the repository.

The core classification cohort encompasses 8,260 radiographs corresponding to all 4,130 unique patients (Table~\ref{tab:data_split}). Stratified splitting was conducted at the patient level using a fixed pseudo-random seed (\texttt{random\_state=42}). To achieve balanced representation of pathology across partitions, each subject was assigned a stratification target defined by the maximum KL grade observed across bilateral knees: $\max(y_{\text{left}}, y_{\text{right}})$. The core cohort was allocated into:
\begin{itemize}[leftmargin=*]
\item \textbf{Training Partition ($\mathcal{P}_{\text{Train}}$):} 2,889 unique patients (5,778 radiographs, 70.0\% of core cohort).
\item \textbf{Validation Partition ($\mathcal{P}_{\text{Val}}$):} 413 unique patients (826 radiographs, 10.0\% of core cohort), reserved exclusively for hyperparameter tuning, learning rate scheduling, and early stopping.
\item \textbf{Primary Held-Out Test Partition ($\mathcal{P}_{\text{Test}}$):} 828 unique patients (1,656 radiographs, 20.0\% of core cohort), retained strictly untouched until terminal evaluation.
\end{itemize}

Pairwise set intersections demonstrate zero subject overlap across all partitions:
\begin{equation}
|\mathcal{P}_{\text{Train}} \cap \mathcal{P}_{\text{Val}}| = |\mathcal{P}_{\text{Train}} \cap \mathcal{P}_{\text{Test}}| = |\mathcal{P}_{\text{Val}} \cap \mathcal{P}_{\text{Test}}| = 0
\label{eq:leakage_free}
\end{equation}

In addition to the core 8,260-image classification set, the Mendeley repository includes a supplementary benchmark partition of 1,526 detector-localized knee joint crops (\texttt{auto\_test}) derived from 811 subjects belonging strictly to the primary held-out test cohort ($\mathcal{P}_{\text{Auto}} \subset \mathcal{P}_{\text{Test}}$) \cite{chen2018dataset}. The benchmark release additionally provides detector-localized knee joint crops, which we use as a supplementary evaluation condition to measure model resilience against localized boundary shifts and field-of-view perturbations under verified zero-training-leakage conditions ($|\mathcal{P}_{\text{Train}} \cap \mathcal{P}_{\text{Auto}}| = 0$).

\begin{table}[t]
\centering
\caption{Patient-level dataset stratification and secondary evaluation subsets.}
\label{tab:data_split}
\scriptsize
\setlength{\tabcolsep}{3.5pt}
\renewcommand{\arraystretch}{1.15}
\resizebox{\columnwidth}{!}{%
\begin{tabular}{lccc}
\toprule
\textbf{Partition} & \textbf{Unique Patients} & \textbf{Radiographs} & \textbf{Core Split (\%)} \\
\midrule
\multicolumn{4}{l}{\textit{Mutually Exclusive Core Classification Cohort:}} \\
Training ($\mathcal{P}_{\text{Train}}$) & 2,889 & 5,778 & 70.0 \\
Validation ($\mathcal{P}_{\text{Val}}$) & 413 & 826 & 10.0 \\
Primary Held-Out Test ($\mathcal{P}_{\text{Test}}$) & 828 & 1,656 & 20.0 \\
\midrule
\textbf{Core Cohort Subtotal} & \textbf{4,130} & \textbf{8,260} & \textbf{100.0} \\
\midrule
\multicolumn{4}{l}{\textit{Supplementary Evaluation Partition:}} \\
Derived Auto-Test Crops$^*$ & 811 & 1,526 & --- \\
\midrule
\textbf{Total Radiograph Archive} & \textbf{4,130} & \textbf{9,786} & --- \\
\bottomrule
\end{tabular}%
}
\vspace{2pt}
\parbox{\columnwidth}{\footnotesize $^*$Auto-Test images represent detector-localized knee joint crops provided within the Mendeley benchmark repository \cite{chen2018dataset}, derived from 811 subjects belonging strictly to the primary held-out test cohort ($\mathcal{P}_{\text{Auto}} \subset \mathcal{P}_{\text{Test}}$). They maintain verified patient-level isolation from training and validation partitions ($|\mathcal{P}_{\text{Train}} \cap \mathcal{P}_{\text{Auto}}| = |\mathcal{P}_{\text{Val}} \cap \mathcal{P}_{\text{Auto}}| = 0$).}
\end{table}

\subsection{Study-Defined Operational Diagnostic Tiers}
\label{subsec:class_remapping}

To align automated prediction with frontline orthopedic triage pathways, the five nominal KL grades were operationally consolidated into three study-defined diagnostic tiers:
\begin{itemize}[leftmargin=*]
\item \textbf{Low/No OA ($\mathcal{Y}_0$, KL~0--1):} Normal joints (KL~0) and doubtful joint-space narrowing with possible marginal osteophytic lipping (KL~1). In clinical practice, this tier reflects non-degenerative or doubtful presentation managed via non-invasive lifestyle guidance and observation \cite{zhang2010eular}.
\item \textbf{Moderate OA ($\mathcal{Y}_1$, KL~2--3):} Definitive radiographic osteoarthritis spanning early osteophytes with preserved joint space (KL~2) to moderate multi-compartment narrowing, sclerosis, and contour deformity (KL~3). This cohort corresponds to conservative non-surgical management, including physical therapy and joint-preservation pharmacotherapy \cite{zhang2010eular}.
\item \textbf{Severe OA ($\mathcal{Y}_2$, KL~4):} Advanced radiographic disease characterized by extensive osteophytosis, marked joint-space loss or obliteration, and severe subchondral sclerosis, identifying cases that may warrant specialist clinical evaluation \cite{zhang2010eular}.
\end{itemize}

Table~\ref{tab:class_dist} reports the class distributions across the primary held-out test partition ($N = 1{,}656$, 828 patients) and the composite multi-partition evaluation cohort ($N = 4{,}008 = 826\ \text{val} + 1{,}656\ \text{test} + 1{,}526\ \text{auto\_test}$). In the primary test partition, Low/No OA comprises 935 samples (56.46\%), Moderate OA comprises 670 samples (40.46\%), and Severe OA comprises 51 samples (3.08\%). In the expanded 4,008-image stress cohort, the distribution spans 2,295 Low/No OA (57.26\%), 1,591 Moderate OA (39.70\%), and 122 Severe OA cases (3.04\%). To mitigate class imbalance during training, batch-wise oversampling was applied via PyTorch \texttt{WeightedRandomSampler} using inverse class frequencies calculated strictly on $\mathcal{P}_{\text{Train}}$, with replacement enabled (\texttt{replacement=True}), drawing $N_{\text{epoch}} = 5{,}778$ balanced samples per epoch prior to data augmentation.

\begin{table}[t]
\centering
\caption{Class distribution across the primary held-out test partition and the composite multi-partition evaluation cohort.}
\label{tab:class_dist}
\scriptsize
\setlength{\tabcolsep}{3pt}
\renewcommand{\arraystretch}{1.15}
\begin{tabular}{lcccc}
\toprule
\multirow{3}{*}{\textbf{Diagnostic Class}} & \multicolumn{2}{c}{\textbf{Held-Out Test}} & \multicolumn{2}{c}{\textbf{Composite Cohort}} \\
 & \multicolumn{2}{c}{\textbf{($N = 1{,}656$)}} & \multicolumn{2}{c}{\textbf{($N = 4{,}008$)}} \\
\cmidrule(lr){2-3} \cmidrule(lr){4-5}
 & \textbf{Samples} & \textbf{\%} & \textbf{Samples} & \textbf{\%} \\
\midrule
Low/No OA (KL 0--1) & 935 & 56.46 & 2{,}295 & 57.26 \\
Moderate OA (KL 2--3) & 670 & 40.46 & 1{,}591 & 39.70 \\
Severe OA (KL 4) & 51 & 3.08 & 122 & 3.04 \\
\midrule
\textbf{Total} & \textbf{1,656} & \textbf{100.00} & \textbf{4,008} & \textbf{100.00} \\
\bottomrule
\end{tabular}
\end{table}

\subsection{Input Preprocessing and Regularization}
\label{subsec:preprocessing}

Input radiographs are standardized to dimension $\mathbf{X} \in \mathbb{R}^{224 \times 224 \times 3}$. Grayscale single-channel radiographs were replicated across three identical RGB channels prior to normalization using ImageNet statistics ($\bm{\mu} = [0.485, 0.456, 0.406]$, $\bm{\sigma} = [0.229, 0.224, 0.225]$). Stochastic training perturbations comprise random horizontal reflection ($p = 0.5$) reflecting bilateral coronal symmetry, minor rotations ($\theta \in [-10^\circ, 10^\circ]$), and brightness/contrast jitter ($\pm 15\%$).

Convex mixup regularization \cite{zhang2018mixup, thulasidasan2019mixup} is applied during training:
\begin{equation}
\tilde{\mathbf{X}} = \lambda \mathbf{X}_i + (1 - \lambda)\mathbf{X}_j, \quad \tilde{\mathbf{y}} = \lambda \mathbf{y}_i^* + (1 - \lambda)\mathbf{y}_j^*
\label{eq:mixup}
\end{equation}
where $\lambda \sim \operatorname{Beta}(\alpha, \alpha)$ with $\alpha = 0.2$. For the ordinal penalty, the mixed ordinal target is computed as $\tilde{r} = \lambda r_i + (1 - \lambda) r_j$, where $r \in \{0, 1, 2\}$ denotes the integer class rank. During validation and testing, strictly deterministic single-sample inference is enforced ($\lambda = 1$).

\subsection{Tri-Paradigm Hybrid Model Architecture}
\label{subsec:architecture}

The proposed architecture (Fig.~\ref{fig:architecture}) constructs a multi-branch representation space by deploying three specialized vision backbones in parallel:

\begin{figure}[!t]
    \centering
    \includegraphics[width=\columnwidth]{fig1.jpeg}
    \caption{End-to-end schematic of the proposed tri-paradigm hybrid vision framework for automated knee osteoarthritis severity grading. A standardized radiograph ($\mathbf{X} \in \mathbb{R}^{224 \times 224 \times 3}$) is concurrently processed under horizontal-flip test-time augmentation through three specialized backbones: ConvNeXt-Base for subchondral bone texture ($D=1024$), Swin-Base Transformer for shifted-window joint compartment context ($D=1024$), and self-supervised DINO ViT-Base for invariant structural priors ($D=768$). The concatenated feature representation ($\mathbf{z}_{\text{fused}} \in \mathbb{R}^{2816}$) is projected through a dual-stage regularized multi-layer perceptron head with batch normalization and dropout, yielding posterior probabilities across three operational diagnostic tiers (Low/No OA KL~0--1, Moderate OA KL~2--3, and Severe OA KL~4) optimized via a distance-aware ordinal composite objective.}
    \label{fig:architecture}
\end{figure}

\subsubsection{Local Texture Operator (ConvNeXt-Base)}
Let $\Phi_{\text{CNN}}(\cdot; \bm{\theta}_C)$ denote the ConvNeXt-Base backbone \cite{liu2022convnext}, parameterized by $\bm{\theta}_C \in \mathbb{R}^{87.6\text{M}}$ (87,566,464 instantiated backbone parameters with classification head removed). ConvNeXt leverages $7 \times 7$ depthwise convolutions with inverted bottleneck channels $[128, 256, 512, 1024]$ across four successive hierarchical stages. The network preserves local translation equivariance, making it exceptionally sensitive to micro-trabecular bone texture and sharp marginal osteophyte spikes. Global average pooling at the terminal stage yields:
\begin{equation}
\mathbf{f}_{\text{CNN}} = \Phi_{\text{CNN}}(\mathbf{X}; \bm{\theta}_C) \in \mathbb{R}^{1024}
\label{eq:f_cnn}
\end{equation}

\subsubsection{Hierarchical Window Attention Operator (Swin Transformer)}
Let $\Phi_{\text{Swin}}(\cdot; \bm{\theta}_S)$ denote the Swin-Base transformer \cite{liu2021swin}, parameterized by $\bm{\theta}_S \in \mathbb{R}^{86.7\text{M}}$ (86,743,224 instantiated backbone parameters). Swin partitions the image into non-overlapping $4 \times 4$ patches and applies shifted-window multi-head self-attention (SW-MSA) within localized $7 \times 7$ token grids, alternating regular and shifted partitions across stages. This captures long-range spatial correlations across the tibiofemoral joint space while maintaining linear computational complexity relative to image area:
\begin{equation}
\mathbf{f}_{\text{Swin}} = \Phi_{\text{Swin}}(\mathbf{X}; \bm{\theta}_S) \in \mathbb{R}^{1024}
\label{eq:f_swin}
\end{equation}

\subsubsection{Self-Supervised Structural Operator (DINO ViT-Base)}
Let $\Phi_{\text{DINO}}(\cdot; \bm{\theta}_D)$ denote the ViT-Base architecture pretrained via DINO self-distillation without human labels \cite{caron2021dino}, parameterized by $\bm{\theta}_D$ with $16 \times 16$ patch embeddings. The instantiated DINO ViT-Base visual backbone utilizes 85.8M feature extractor parameters (85,798,656 with pretraining projection head removed), maintained completely frozen during training to preserve self-supervised structural priors. Rather than optimizing cross-entropy on natural categories, DINO's momentum-teacher centering objective induces emergent attention toward object contours, articular boundaries, and structural landmarks. In our framework, $\bm{\theta}_D$ is maintained completely frozen during training to preserve these universal structural priors:
\begin{equation}
\mathbf{f}_{\text{DINO}} = \Phi_{\text{DINO}}(\mathbf{X}; \bm{\theta}_D^*) \in \mathbb{R}^{768}
\label{eq:f_dino}
\end{equation}

\subsubsection{Late Manifold Fusion and Multi-Layer Projection Head}
The three visual representations are concatenated into a unified latent feature manifold $\mathbf{z}_{\text{fused}} \in \mathbb{R}^{D}$ with total dimensionality $D = 1024 + 1024 + 768 = 2816$:
\begin{equation}
\mathbf{z}_{\text{fused}} = \left[ \mathbf{f}_{\text{CNN}}^\top, \, \mathbf{f}_{\text{Swin}}^\top, \, \mathbf{f}_{\text{DINO}}^\top \right]^\top \in \mathbb{R}^{2816}
\label{eq:concat}
\end{equation}

The classification projection head $g_{\bm{\psi}}: \mathbb{R}^{2816} \to \mathbb{R}^3$ comprises two hidden fully connected layers interleaved with Batch Normalization ($\operatorname{BN}$), Rectified Linear Units ($\operatorname{ReLU}$), and Bernoulli dropout operators $\mathbf{m}_l \sim \operatorname{Bernoulli}(1 - p_l)$:
\begin{align}
\mathbf{h}_1 &= \mathbf{m}_1 \odot \operatorname{ReLU}\left(\operatorname{BN}\left(\mathbf{W}_1 \mathbf{z}_{\text{fused}} + \mathbf{b}_1\right)\right), \quad \mathbf{h}_1 \in \mathbb{R}^{1024} \label{eq:head1}\\
\mathbf{h}_2 &= \mathbf{m}_2 \odot \operatorname{ReLU}\left(\operatorname{BN}\left(\mathbf{W}_2 \mathbf{h}_1 + \mathbf{b}_2\right)\right), \quad \mathbf{h}_2 \in \mathbb{R}^{512} \label{eq:head2}\\
\mathbf{s}(\mathbf{X}) &= \mathbf{W}_3 \mathbf{h}_2 + \mathbf{b}_3, \quad \mathbf{s}(\mathbf{X}) \in \mathbb{R}^{3} \label{eq:logits}
\end{align}
where $p_1 = 0.5$, $p_2 = 0.4$, $\mathbf{W}_1 \in \mathbb{R}^{1024 \times 2816}$, $\mathbf{W}_2 \in \mathbb{R}^{512 \times 1024}$, and $\mathbf{W}_3 \in \mathbb{R}^{3 \times 512}$. Accounting for the fully connected linear transformation tensors, bias vectors, and affine Batch Normalization parameters ($2 \times 1024 + 2 \times 512 = 3{,}072$), the dual-stage projection head comprises exactly 3,414,019 parameters ($\approx 3.41\text{M}$). Posterior severity probabilities $\mathbf{p} \in \Delta^2$ on the probability simplex are obtained via softmax:
\begin{equation}
p_k(\mathbf{X}) = \frac{\exp(s_k(\mathbf{X}))}{\sum_{j=0}^{K-1} \exp(s_j(\mathbf{X}))}, \quad k \in \{0, 1, 2\}
\label{eq:softmax}
\end{equation}

\subsection{Formal Ordinal-Aware Loss Formulation}
\label{subsec:loss}

Standard categorical cross-entropy treats label discrepancies nominally, assigning equal loss regardless of whether a Grade~4 radiograph is misclassified as Grade~3 or Grade~0. To enforce clinical safety, we formulate an ordinal risk objective that incorporates continuous metric distance in label space.

Let $\mathbf{y}^* \in \Delta^{K-1}$ denote the ground-truth target vector under label smoothing regularizer $\epsilon = 0.05$:
\begin{equation}
y_k^* = (1 - \epsilon)\mathbb{I}(y = k) + \frac{\epsilon}{K}, \quad k \in \{0, 1, \dots, K-1\}
\label{eq:smoothed_label}
\end{equation}

The cross-entropy component is given by:
\begin{equation}
\mathcal{L}_{\text{CE}}(\mathbf{p}, \mathbf{y}^*) = -\sum_{k=0}^{K-1} y_k^* \log p_k
\label{eq:ce_loss}
\end{equation}

We define the expected ordinal severity rank $\hat{y} \in [0, K-1]$ under the model's posterior distribution:
\begin{equation}
\hat{y}(\mathbf{X}) = \mathbb{E}_{k \sim \mathbf{p}}[k] = \sum_{k=0}^{K-1} k \cdot p_k(\mathbf{X})
\label{eq:expected_rank}
\end{equation}

The ordinal distance penalty is formulated as the expected absolute error between continuous prediction $\hat{y}$ and true ordinal grade $y \in \{0, 1, \dots, K-1\}$:
\begin{equation}
\mathcal{D}_{\text{ord}}(\mathbf{p}, y) = |\hat{y}(\mathbf{X}) - y| = \left| \sum_{k=0}^{K-1} k \cdot p_k(\mathbf{X}) - y \right|
\label{eq:mae_loss}
\end{equation}

The complete empirical risk functional across a mini-batch of size $B$ is:
\begin{equation}
\mathcal{L}_{\text{total}} = \frac{1}{B}\sum_{i=1}^B \left[ \mathcal{L}_{\text{CE}}(\mathbf{p}^{(i)}, \tilde{\mathbf{y}}_i) + \lambda_{\text{ord}} \cdot \mathcal{D}_{\text{ord}}(\mathbf{p}^{(i)}, \tilde{r}_i) \right] + \frac{\gamma}{2}\|\bm{\Theta}\|_2^2
\label{eq:total_loss}
\end{equation}
where $\lambda_{\text{ord}} = 0.5$ governs the ordinal penalty weight, $\gamma = 10^{-4}$ represents weight decay, and $\bm{\Theta} = \{\bm{\theta}_C, \bm{\theta}_S, \bm{\psi}\}$ denotes trainable parameters. During training with convex Mixup data augmentation ($\lambda < 1$), the cross-entropy term optimizes against the blended soft target distribution $\tilde{\mathbf{y}}_i$ \eqref{eq:mixup}, while the ordinal penalty minimizes metric error against the corresponding blended ordinal rank $\tilde{r}_i = \lambda r_i + (1 - \lambda)r_j$: $\mathcal{D}_{\text{ord}}(\mathbf{p}^{(i)}, \tilde{r}_i) = |\hat{y}(\mathbf{X}^{(i)}) - \tilde{r}_i|$. During deterministic validation and testing ($\lambda = 1$), $\tilde{\mathbf{y}}_i = \mathbf{y}_i^*$ and $\tilde{r}_i = y_i$. Under this formulation, misclassifying a Severe ($y=2$) radiograph as Low/No OA ($k=0$) incurs a metric distance penalty twice as large as classifying it as Moderate ($k=1$), imposing a distance-proportional mathematical penalty that suppresses catastrophic clinical grade omissions and steers optimization away from non-adjacent prediction errors.

To rigorously characterize the optimization dynamics induced by the ordinal penalty, we derive the analytical gradient of $\mathcal{L}_{\text{total}}$ with respect to the network logits $\mathbf{s}(\mathbf{X}) \in \mathbb{R}^K$. Using the chain rule, the gradient decomposes into cross-entropy and ordinal penalty components:
\begin{equation}
\frac{\partial \mathcal{L}_{\text{total}}}{\partial s_k} = \frac{\partial \mathcal{L}_{\text{CE}}}{\partial s_k} + \lambda_{\text{ord}} \frac{\partial \mathcal{D}_{\text{ord}}}{\partial s_k}
\label{eq:gradient_decomp}
\end{equation}

For the label-smoothed cross-entropy loss, standard differentiation yields:
\begin{equation}
\frac{\partial \mathcal{L}_{\text{CE}}}{\partial s_k} = p_k - y_k^*
\label{eq:grad_ce}
\end{equation}

For the ordinal penalty $\mathcal{D}_{\text{ord}}(\mathbf{p}, y) = |\hat{y}(\mathbf{X}) - y|$, applying the subgradient operator gives:
\begin{equation}
\frac{\partial \mathcal{D}_{\text{ord}}}{\partial s_k} = \operatorname{sgn}(\hat{y}(\mathbf{X}) - y) \cdot \frac{\partial \hat{y}(\mathbf{X})}{\partial s_k}
\label{eq:grad_ord_chain}
\end{equation}
where $\operatorname{sgn}(\cdot)$ is the signum function. Expanding the derivative of the continuous expectation $\hat{y}(\mathbf{X}) = \sum_{j=0}^{K-1} j \cdot p_j$ with respect to logit $s_k$ via the Jacobian of the softmax mapping $\frac{\partial p_j}{\partial s_k} = p_k(\delta_{jk} - p_j)$:
\begin{align}
\frac{\partial \hat{y}(\mathbf{X})}{\partial s_k} &= \sum_{j=0}^{K-1} j \cdot \frac{\partial p_j}{\partial s_k} = \sum_{j=0}^{K-1} j \cdot p_k (\delta_{jk} - p_j) \nonumber \\
&= k \cdot p_k - p_k \sum_{j=0}^{K-1} j \cdot p_j = p_k \left( k - \hat{y}(\mathbf{X}) \right)
\label{eq:grad_expected_y}
\end{align}

Substituting \eqref{eq:grad_ce} and \eqref{eq:grad_expected_y} into \eqref{eq:gradient_decomp} yields the complete analytical logit gradient:
\begin{equation}
\frac{\partial \mathcal{L}_{\text{total}}}{\partial s_k} = (p_k - y_k^*) + \lambda_{\text{ord}} \cdot \operatorname{sgn}(\hat{y}(\mathbf{X}) - y) \cdot p_k \left( k - \hat{y}(\mathbf{X}) \right)
\label{eq:full_gradient}
\end{equation}

Equation~\eqref{eq:full_gradient} reveals the underlying mechanism of catastrophic error suppression: when an advanced radiograph ($y=2$, Severe) is severely mispredicted toward the healthy boundary ($\hat{y} \approx 0 < y$), the directional multiplier becomes $\operatorname{sgn}(\hat{y} - y) = -1$. For the true class logit ($k=2$), the bracketed term evaluates to $(2 - \hat{y}) > 0$, delivering an amplifying negative gradient $-\lambda_{\text{ord}} p_2 (2 - \hat{y})$ that drives the logit $s_2$ upward with force proportional to distance. Conversely, for the non-adjacent false class ($k=0$), $k - \hat{y} \approx 0$, suppressing benign updates. Consequently, the gradient field automatically exerts an asymmetric repulsive force against extreme non-adjacent transitions, mathematically underpinning the zero-reversal empirical finding.

\subsection{Two-Stage Progressive Fine-Tuning}
\label{subsec:training}

To prevent catastrophic disruption of ImageNet and self-supervised representations by randomly initialized head weights, optimization proceeds in two sequential phases:

\begin{itemize}[leftmargin=*]
\item \textbf{Phase 1 (Warmup):} All backbone weights $\{\bm{\theta}_C, \bm{\theta}_S, \bm{\theta}_D\}$ remain frozen. Only projection head parameters $\bm{\psi}$ are updated for 5 epochs using AdamW \cite{loshchilov2019adamw} with learning rate $\eta_{\text{warmup}} = 10^{-3}$.
\item \textbf{Phase 2 (Selective Fine-Tuning):} The terminal stage of ConvNeXt (\texttt{stages.3}) and the final layer of Swin Transformer (\texttt{layers.3}) are unfrozen alongside $\bm{\psi}$, while earlier stages and DINO ViT remain frozen. Training proceeds for 20 epochs with base learning rate $\eta_{\text{base}} = 5 \times 10^{-5}$ under cosine annealing \cite{loshchilov2017sgdr}:
\begin{equation}
\eta_t = \eta_{\min} + \frac{1}{2}(\eta_{\text{base}} - \eta_{\min})\left(1 + \cos\left(\frac{t \cdot \pi}{T}\right)\right)
\label{eq:cosine}
\end{equation}
where $\eta_{\min} = 10^{-6}$ and $T = 20$. Mixed precision (FP16) via PyTorch \texttt{GradScaler} was utilized throughout.
\item \textbf{Pretrained Checkpoints and Standardization:} Backbones were instantiated via the \texttt{timm} library (v0.9.16) using official checkpoints: \texttt{convnext\_base.fb\_in22k\_ft\_in1k} for ConvNeXt-Base, \texttt{swin\_base\_patch4\_window7\_224} for Swin Transformer, and \texttt{vit\_base\_patch16\_224.dino} (DINO v1) for self-supervised ViT-Base. Standalone baselines and the hybrid model were trained under identical patient-level partitions, data augmentations, loss formulations, AdamW optimizers, and learning rate schedules.
\item \textbf{Reproducibility Protocol and Checkpoint Selection:} Global random seeds were locked across Python \texttt{random}, NumPy, PyTorch, and CUDA backend (\texttt{torch.backends.cudnn.deterministic = True}) using seed 42. Terminal model weights were selected strictly based on the highest validation Quadratic Weighted Kappa (QWK) on $\mathcal{P}_{\text{Val}}$. To assess stability against partition stochasticity without contaminating the evaluation benchmark, five-fold cross-validation was conducted strictly within the development cohort ($\mathcal{P}_{\text{Train}} \cup \mathcal{P}_{\text{Val}}$, $N = 6{,}604$ radiographs across 3,302 patients) using stratified patient-level folds. Cross-validation yielded consistent out-of-fold validation performance (mean validation accuracy $84.82 \pm 0.45\%$, mean QWK $0.743 \pm 0.012$ under TTA), demonstrating architectural resilience against split stochasticity while maintaining the primary held-out test partition ($\mathcal{P}_{\text{Test}}$, $N = 1{,}656$, 828 patients) completely untouched until terminal evaluation. For each fold, a distinct model was trained from scratch, all hyperparameter tuning and model selection were performed strictly within that iteration's training partition, and validation metrics were evaluated on the corresponding out-of-fold split.
\end{itemize}

\subsection{Dual-Pass Test-Time Augmentation and Ensemble Consensus}
\label{subsec:tta}

At inference, prediction variance along ambiguous decision boundaries is reduced via horizontal-flip test-time augmentation (TTA) \cite{shanmugam2021tta}. Let $\mathbf{R}_{\text{horiz}}$ denote the horizontal reflection operator:
\begin{equation}
\mathbf{p}_{\text{TTA}}(\mathbf{X}) = \frac{1}{2}\left[ \mathbf{p}(\mathbf{X}) + \mathbf{p}(\mathbf{R}_{\text{horiz}}\mathbf{X}) \right]
\label{eq:tta}
\end{equation}

The soft-voting ensemble prediction across the four architectural configurations $\mathcal{M} = \{\text{CNN}, \text{Swin}, \text{DINO}, \text{Hybrid}\}$ is derived as a weighted convex combination:
\begin{equation}
\mathbf{p}_{\text{ensemble}}(\mathbf{X}) = \sum_{m \in \mathcal{M}} w_m \mathbf{p}_{\text{TTA}}^{(m)}(\mathbf{X})
\label{eq:ensemble}
\end{equation}
where weights $\mathbf{w} = [0.10, 0.20, 0.10, 0.60]$ were determined via grid search with a step size of $\Delta w = 0.05$ over the probability simplex strictly on the validation partition ($N = 826$ images, 413 patients) by maximizing validation Quadratic Weighted Kappa ($\kappa_{\text{QWK}}$), granting dominant influence to the tri-paradigm hybrid model while maintaining the primary held-out test partition untouched for unbiased evaluation.

\subsection{Multi-Method Explainability Suite}
\label{subsec:xai}

To audit model fidelity against clinical biomarkers in alignment with radiology interpretability principles \cite{reyes2020interpretability}, seven complementary attribution methods are implemented. The attribution analysis focuses on input-level perturbations and terminal ConvNeXt feature-space activations; dedicated branch-specific attribution decoupling for Swin and DINO remains an avenue for future work:
\begin{enumerate}[leftmargin=*, label=(\roman*)]
\item \textbf{Grad-CAM} \cite{selvaraju2017gradcam}: Gradient-weighted channel activations in the terminal ConvNeXt stage.
\item \textbf{Grad-CAM++} \cite{chattopadhay2018gradcampp}: Higher-order partial derivative weighting for multi-focal osteophyte localization.
\item \textbf{Layer-CAM} \cite{jiang2021layercam}: Spatially granular class activation preserving fine-grained trabecular edge detail.
\item \textbf{Score-CAM} \cite{wang2020scorecam}: Gradient-free perturbation attribution evaluating channel-wise forward confidence scores.
\item \textbf{Integrated Gradients} \cite{sundararajan2017axiomatic}: Axiomatically complete path-integral attributions accumulated from a blank baseline:
\begin{equation}
\operatorname{IG}_i(\mathbf{X}) = (X_i - X'_i) \times \int_0^1 \frac{\partial F(\mathbf{X}' + \alpha(\mathbf{X} - \mathbf{X}'))}{\partial X_i} d\alpha
\label{eq:ig}
\end{equation}
\item \textbf{Vanilla Saliency}: Direct pixel sensitivity magnitudes $|\partial s_k / \partial \mathbf{X}|$.
\item \textbf{LIME Superpixel Explanations} \cite{ribeiro2016lime}: Model-agnostic local surrogate modeling. Input images are segmented via SLIC into superpixels; 300 perturbed instances are sampled and fitted via Ridge regression to identify positively and negatively contributing anatomical compartments.
\end{enumerate}


% ==============================================================================
% 4. RESULTS
% ==============================================================================
\section{Results}
\label{sec:results}

\subsection{Overall Model Performance on Primary Held-Out Test Partition}
\label{subsec:overall_performance}

Table~\ref{tab:overall_performance} provides a rigorous comparison between single-pass (static inference) and horizontal-flip test-time augmentation (TTA) across all evaluated architectures on the strictly isolated primary held-out test partition ($N = 1{,}656$ radiographs from 828 unique patients). All performance metrics are accompanied by 95\% confidence intervals (CIs) estimated via patient-level bootstrap resampling ($B = 1{,}000$ iterations) to account for within-patient correlation between bilateral knees.

\begin{table*}[t]
\centering
\caption{Model performance on the primary patient-level held-out test partition ($N = 1{,}656$ radiographs, 828 unique patients) under single-pass static inference versus horizontal-flip test-time augmentation (TTA). Metric uncertainty is reported as 95\% confidence intervals [in brackets] computed via 1,000 patient-level bootstrap iterations. Accuracy, F1 scores, Sensitivity, and Specificity are expressed as percentages; QWK and Macro AUC are decimal scores.}
\label{tab:overall_performance}
\scriptsize
\setlength{\tabcolsep}{3pt}
\renewcommand{\arraystretch}{1.1}
\resizebox{\textwidth}{!}{%
\begin{tabular}{llccccccc}
\toprule
\textbf{Model Architecture} & \textbf{Inference Mode} & \textbf{Accuracy (\%)} & \textbf{Macro F1 (\%)} & \textbf{Weighted F1 (\%)} & \textbf{Sensitivity (\%)} & \textbf{Specificity (\%)} & \textbf{QWK} & \textbf{Macro AUC} \\
\midrule
\multirow{2}{*}{ConvNeXt-Base} 
& Static & 83.78 [81.98, 85.58] & 80.52 [78.41, 82.63] & 83.66 [81.85, 85.47] & 84.22 [81.95, 86.49] & 89.12 [87.45, 90.79] & 0.725 [0.698, 0.752] & 0.941 [0.929, 0.953] \\
& Horiz-Flip TTA & 84.90 [83.15, 86.65] & 81.65 [79.60, 83.70] & 84.70 [82.95, 86.45] & 85.34 [83.12, 87.56] & 89.85 [88.23, 91.47] & 0.742 [0.716, 0.768] & 0.950 [0.939, 0.961] \\
\midrule
\multirow{2}{*}{Swin Transformer} 
& Static & 80.29 [78.36, 82.22] & 77.41 [75.18, 79.64] & 79.85 [77.92, 81.78] & 84.13 [81.82, 86.44] & 86.46 [84.65, 88.27] & 0.672 [0.643, 0.701] & 0.922 [0.908, 0.936] \\
& Horiz-Flip TTA & 81.76 [79.88, 83.64] & 78.82 [76.65, 80.99] & 81.40 [79.52, 83.28] & 85.40 [83.15, 87.65] & 87.52 [85.78, 89.26] & 0.695 [0.667, 0.723] & 0.934 [0.921, 0.947] \\
\midrule
\multirow{2}{*}{DINO ViT} 
& Static & 83.09 [81.27, 84.91] & 83.69 [81.72, 85.66] & 83.02 [81.18, 84.86] & 86.91 [84.82, 89.00] & 88.73 [87.05, 90.41] & 0.716 [0.689, 0.743] & 0.940 [0.928, 0.952] \\
& Horiz-Flip TTA & 84.42 [82.66, 86.18] & 84.85 [82.95, 86.75] & 84.38 [82.60, 86.16] & 88.10 [86.08, 90.12] & 89.60 [87.98, 91.22] & 0.736 [0.710, 0.762] & 0.949 [0.938, 0.960] \\
\midrule
\multirow{2}{*}{\textbf{Hybrid Fusion}} 
& Static & 83.44 [81.65, 85.23] & 83.38 [81.42, 85.34] & 83.21 [81.40, 85.02] & 86.73 [84.62, 88.84] & 88.61 [86.92, 90.30] & 0.721 [0.694, 0.748] & 0.943 [0.931, 0.955] \\
& \textbf{Horiz-Flip TTA} & \textbf{85.08 [83.35, 86.81]} & \textbf{84.95 [83.08, 86.82]} & \textbf{84.88 [83.15, 86.61]} & \textbf{88.42 [86.45, 90.39]} & \textbf{89.92 [88.32, 91.52]} & \textbf{0.748 [0.722, 0.774]} & \textbf{0.954 [0.943, 0.965]} \\
\midrule
\multirow{2}{*}{\textbf{Weighted Ensemble}} 
& Static & 84.19 [82.42, 85.96] & 83.19 [81.20, 85.18] & 83.96 [82.18, 85.74] & 87.20 [85.12, 89.28] & 89.12 [87.46, 90.78] & 0.735 [0.708, 0.762] & 0.948 [0.937, 0.959] \\
& \textbf{Horiz-Flip TTA} & \textbf{85.63 [83.93, 87.33]} & \textbf{85.20 [83.32, 87.08]} & \textbf{85.45 [83.74, 87.16]} & \textbf{88.85 [86.90, 90.80]} & \textbf{90.25 [88.68, 91.82]} & \textbf{0.758 [0.733, 0.783]} & \textbf{0.958 [0.948, 0.968]} \\
\bottomrule
\end{tabular}%
}
\end{table*}

Under static single-pass inference, the tri-paradigm Hybrid model attains an overall accuracy of 83.44\% (95\% CI: [81.65\%, 85.23\%]), a Quadratic Weighted Kappa (QWK) of 0.721 [0.694, 0.748], and a macro ROC-AUC of 0.943 [0.931, 0.955]. Standalone ConvNeXt-Base achieves 83.78\% accuracy, DINO ViT yields 83.09\%, and the Soft-Voting Ensemble reaches 84.19\% accuracy. Applying horizontal-flip TTA on this exact same 1,656-image cohort produces a consistent performance gain across all architectures, advancing Hybrid Fusion accuracy by +1.64 percentage points (pp) to 85.08\% (QWK: 0.748, AUC: 0.954) and Weighted Ensemble accuracy by +1.44 pp to 85.63\% (QWK: 0.758, AUC: 0.958). This controlled comparison isolates the genuine algorithmic benefit of test-time coronal reflection from variations in sample size or cohort composition.

\subsection{Multi-Partition Stress-Testing on the Composite Cohort}
\label{subsec:tta_impact}

To evaluate model resilience against multi-source imaging variations, localized detector cropping noise, and broader sample diversity, Table~\ref{tab:tta_impact} reports performance under GPU dual-pass TTA on the composite multi-partition evaluation cohort ($N = 4{,}008$ radiographs, comprising 826 validation, 1,656 primary held-out test, and 1,526 detector-derived auto-test images). Because validation images were utilized for ensemble weight calibration and auto-test images represent automated YOLO crops from test subjects, this 4,008-image cohort is presented strictly as an expanded stress-testing evaluation rather than an independent unselected benchmark.

\begin{table*}[t]
\centering
\caption{Model performance under GPU horizontal-flip TTA on the composite multi-partition evaluation cohort ($N = 4{,}008$ radiographs, comprising validation, primary held-out test, and detector-derived auto-test partitions). Metric uncertainty is reported as 95\% confidence intervals [in brackets] estimated using patient-level cluster bootstrap resampling ($B = 1{,}000$ iterations), jointly resampling all original and derived detector crops belonging to each unique patient identifier. All values are percentages except QWK and ROC-AUC.}
\label{tab:tta_impact}
\scriptsize
\setlength{\tabcolsep}{3.5pt}
\renewcommand{\arraystretch}{1.1}
\resizebox{\textwidth}{!}{%
\begin{tabular}{lcccccc}
\toprule
\textbf{Model Architecture} & \textbf{Accuracy (\%)} & \textbf{Precision (\%)} & \textbf{Recall (\%)} & \textbf{Macro F1 (\%)} & \textbf{QWK} & \textbf{Macro AUC} \\
\midrule
ConvNeXt-Base & 87.52 [86.48, 88.56] & 88.07 [86.85, 89.29] & 90.27 [89.15, 91.39] & 89.05 [87.95, 90.15] & 0.800 [0.785, 0.815] & 0.966 [0.959, 0.973] \\
Swin Transformer & 84.26 [83.12, 85.40] & 83.00 [81.65, 84.35] & 87.28 [86.05, 88.51] & 84.35 [83.12, 85.58] & 0.749 [0.732, 0.766] & 0.948 [0.939, 0.957] \\
DINO ViT & 90.22 [89.28, 91.16] & 92.08 [91.15, 93.01] & 92.98 [92.05, 93.91] & 92.52 [91.65, 93.39] & 0.842 [0.829, 0.855] & 0.974 [0.968, 0.980] \\
\textbf{Hybrid Fusion} & \textbf{90.92 [90.02, 91.82]} & \textbf{92.39 [91.48, 93.30]} & \textbf{93.17 [92.26, 94.08]} & \textbf{92.70 [91.85, 93.55]} & \textbf{0.853 [0.841, 0.865]} & \textbf{0.980 [0.975, 0.985]} \\
Soft-Voting Ensemble & 90.54 [89.62, 91.46] & 92.02 [91.08, 92.96] & 92.80 [91.86, 93.74] & 92.28 [91.40, 93.16] & 0.847 [0.835, 0.859] & 0.980 [0.975, 0.985] \\
\bottomrule
\end{tabular}%
}
\end{table*}

On this expanded cohort, Hybrid Fusion achieved the highest observed accuracy (90.92\%, 95\% CI: [90.02\%, 91.82\%]), compared with 90.54\% for the Soft-Voting Ensemble and 90.22\% for DINO ViT. Standalone ConvNeXt-Base achieves 87.52\%, while Swin Transformer attains 84.26\%. The QWK of 0.853 reflects substantial ordinal concordance according to the Landis--Koch scale \cite{landis1977measurement} across the three diagnostic tiers, while the macro ROC-AUC of 0.980 reflects robust discriminative capacity across multi-partition inputs.

\subsection{Class-Wise Performance Analysis}
\label{subsec:classwise}

Table~\ref{tab:classwise} presents the class-stratified precision, recall, and F1-score for the Soft-Voting Ensemble under horizontal-flip TTA on the composite evaluation cohort ($N = 4{,}008$).

\begin{table}[t]
\centering
\caption{Per-class classification report for the Soft-Voting Ensemble with horizontal-flip TTA on the composite evaluation cohort ($N = 4{,}008$).}
\label{tab:classwise}
\scriptsize
\setlength{\tabcolsep}{5pt}
\renewcommand{\arraystretch}{1.1}
\begin{tabular}{lcccc}
\toprule
\textbf{Diagnostic Tier} & \textbf{Precision} & \textbf{Recall} & \textbf{F1-Score} & \textbf{Support} \\
\midrule
Low/No OA (KL 0--1) & 0.8945 & 0.9499 & 0.9214 & 2{,}295 \\
Moderate OA (KL 2--3) & 0.9202 & 0.8341 & 0.8750 & 1{,}591 \\
\textbf{Severe OA (KL 4)} & \textbf{0.9457} & \textbf{1.0000} & \textbf{0.9721} & \textbf{122} \\
\midrule
Accuracy & \multicolumn{3}{c}{0.9054} & 4{,}008 \\
Macro Average & 0.9202 & 0.9280 & 0.9228 & 4{,}008 \\
Weighted Average & 0.9063 & 0.9054 & 0.9045 & 4{,}008 \\
\bottomrule
\end{tabular}
\end{table}

Key observations from the per-tier analysis include:
\begin{enumerate}[leftmargin=*]
\item \textbf{Primary Test and Stress-Cohort Severe Discrimination:} On the strictly isolated primary held-out test partition ($N = 1{,}656$, 828 unique patients), Severe OA recall reached 100.00\% ($51/51$ correct detections, 95\% CI: [94.1\%, 100.0\%]), with zero non-adjacent omissions into Low/No OA ($0/51$) and precision of 94.44\% ($51/54$, F1-score: 97.14\%). On the expanded 4,008-image composite stress cohort, this boundary containment generalized across all 122 Severe instances ($122/122$, 100.00\% empirical sensitivity, 95\% CI: [97.0\%, 100.0\%], precision 94.57\%, F1-score 97.21\%). No non-adjacent Severe-to-Low/No OA errors were observed in the evaluated configurations across either the independent held-out test partition or the detector-perturbed derived evaluation cohort. The controlled loss ablation in Table~\ref{tab:loss_ablation} provides evidence that the distance-aware ordinal formulation was associated with fewer observed non-adjacent grade reversals in the controlled ablation, where omitting the ordinal penalty ($\lambda_{\text{ord}} = 0$) permitted non-adjacent misclassifications into Low/No OA.
\item \textbf{Error Distribution along Borderline Boundaries:} For the Soft-Voting Ensemble, overall misclassifications in the Moderate cohort totaled 16.59\% (264/1,591), comprising 257 cases assigned to Low/No OA (16.15\%) and 7 cases assigned to Severe OA (0.44\%). These transitional disagreements correspond to the clinically documented difficulty of distinguishing doubtful osteophytic lipping from definitive early OA (KL~1 vs.\ KL~2).
\end{enumerate}

\subsection{Confusion Matrix and ROC-AUC Analysis}
\label{subsec:confusion_roc}

\begin{figure*}[!t]
    \centering
    \includegraphics[width=0.95\textwidth]{fig2.jpeg}
    \caption{Empirical confusion matrices across all five evaluated model configurations on the composite evaluation cohort ($N = 4{,}008$ radiographs across 2,295 Low/No OA, 1,591 Moderate OA, and 122 Severe OA instances): (a) standalone ConvNeXt-Base (87.52\% accuracy with TTA), (b) standalone Swin-Base (84.26\%), (c) standalone DINO ViT-Base (90.22\%), (d) proposed Hybrid Fusion model (90.92\%), and (e) Soft-Voting Ensemble (90.54\%). Non-zero off-diagonal entries are strictly restricted to adjacent diagnostic boundaries, achieving zero extreme grade reversals (0/122 Severe cases misclassified as Low/No OA) and 100\% empirical sensitivity for advanced disease across DINO ViT, Hybrid Fusion, and Ensemble configurations.}
    \label{fig:confusion}
\end{figure*}

Figure~\ref{fig:confusion} details the empirical confusion matrices across all five model configurations on the composite $N = 4{,}008$ evaluation cohort. Across all five pipelines, classification dispersion is strictly bounded to immediately adjacent diagnostic tiers. On the Severe cohort ($N = 122$), standalone ConvNeXt-Base correctly identifies 120/122 instances (98.36\%), misclassifying only 2 cases into the adjacent Moderate tier and zero into the Low/No OA tier. Standalone Swin-Base mirrors this distribution with 120/122 Severe instances correctly detected. Standalone DINO ViT-Base, the proposed Hybrid Fusion model, and the Soft-Voting Ensemble all achieve 100.00\% empirical sensitivity (122/122 correct detections) on the Severe cohort, completely eliminating false negatives for joint-replacement referral.

For the Low/No OA tier ($N = 2{,}295$), the proposed Hybrid Fusion model attains a diagnostic recall of 94.60\%, correctly identifying 2,171 instances, with all 124 residual errors confined to the borderline Moderate tier and zero into Severe OA. In the Moderate cohort ($N = 1{,}591$), Hybrid Fusion correctly identifies 1,351 cases (84.92\% recall), with 234 misclassified as Low/No OA (14.71\%) and 6 misclassified as Severe (0.38\%). Both error directions remain strictly adjacent to the ground-truth tier, yielding a total Moderate error rate of 15.08\% (240/1,591). The Soft-Voting Ensemble demonstrates comparable consistency, correctly identifying 2,165 Low/No OA cases (94.34\% recall), 1,327 Moderate cases (83.41\% recall; with 257 assigned to Low/No OA [16.15\%] and 7 assigned to Severe [0.44\%]), and 122 Severe cases (100.00\% recall). This empirical distribution supports the contribution of the distance-aware ordinal penalty to suppressing non-adjacent prediction errors, concentrating residual uncertainty entirely along the clinically ambiguous KL~1 versus KL~2 transition boundary.

\begin{figure*}[!b]
    \centering
    \includegraphics[width=0.95\textwidth]{fig3.jpeg}
    \caption{Multi-class one-versus-rest Receiver Operating Characteristic (ROC) curves across all five model configurations under GPU horizontal-flip test-time augmentation ($N = 4{,}008$). Individual area under the curve (AUC) metrics are displayed for Low/No OA (KL~0--1), Moderate OA (KL~2--3), and Severe OA (KL~4) cohorts alongside the overall Macro-Averaged AUC. Both the proposed Hybrid Fusion model and the Soft-Voting Ensemble attain a macro AUC of 0.980, with Severe osteoarthritis reaching an AUC of 1.000.}
    \label{fig:roc}
\end{figure*}

Figure~\ref{fig:roc} depicts the corresponding one-versus-rest ROC curves evaluated under GPU horizontal-flip TTA. Standalone ConvNeXt-Base achieves AUCs of 0.957 for Low/No OA, 0.942 for Moderate OA, and 0.999 for Severe OA (macro AUC = 0.966). Standalone Swin-Base yields AUCs of 0.939, 0.916, and 0.989 (macro AUC = 0.948). Standalone DINO ViT-Base exhibits substantial discriminative power with AUCs of 0.966 for Low/No OA, 0.956 for Moderate OA, and 1.000 for Severe OA (macro AUC = 0.974). Both the proposed Hybrid Fusion model and the Soft-Voting Ensemble achieve the peak macro AUC of 0.980, supported by individual class AUCs of 0.970 for Low/No OA, 0.968 for Moderate OA, and 1.000 for Severe OA.

\subsection{Statistical Significance and Multiple-Testing Correction}
\label{subsec:stats_sig}

Pairwise statistical differences between individual backbones and the Ensemble were evaluated via McNemar's test with Edwards' continuity correction on the primary held-out test partition ($N = 1{,}656$):
\begin{equation}
\chi_{\text{corrected}}^2 = \frac{\left( |n_{10} - n_{01}| - 1 \right)^2}{n_{10} + n_{01}} \sim \chi^2(1)
\label{eq:mcnemar_stat}
\end{equation}
where $n_{10}$ denotes samples correct under the baseline model but failed by the ensemble, and $n_{01}$ denotes the converse. To control the family-wise error rate across $M = 4$ comparisons, Bonferroni-Holm step-down adjustment was applied (Table~\ref{tab:mcnemar}). Because bilateral radiographs from the same patient share anatomical habitus, image-level McNemar testing serves as an approximate assessment of classification divergence; patient-stratified bootstrap resampling yields consistent conclusions.

\begin{table}[!t]
\centering
\caption{Pairwise McNemar's statistical significance tests comparing individual backbone architectures against the Fused Ensemble on the primary held-out test partition ($N=1{,}656$, $\text{df} = 1$, with Edwards' continuity correction and Bonferroni-Holm adjustment).}
\label{tab:mcnemar}
\scriptsize
\setlength{\tabcolsep}{3pt}
\renewcommand{\arraystretch}{1.1}
\begin{tabularx}{\columnwidth}{@{}p{0.32\columnwidth}ccccX@{}}
\toprule
\textbf{Comparison} & $\bm{n_{10}/n_{01}}$ & $\bm{\chi^2}$ & \textbf{Raw} $\bm{p}$ & $\bm{\alpha_{\text{Holm}}}$ & \textbf{Inference} \\
\midrule
Swin vs. Ensemble & $21/78$ & $31.525$ & $<10^{-4}$ & $0.0125$ & Significant ($p < 0.001$) \\
DINO vs. Ensemble & $22/38$ & $3.750$ & $0.0528$ & $0.0167$ & Not significant \\
Hybrid vs. Ensemble & $15/26$ & $2.439$ & $0.1184$ & $0.0250$ & Not significant \\
ConvNeXt vs. Ensemble & $18/24$ & $0.595$ & $0.4404$ & $0.0500$ & Not significant \\
\bottomrule
\end{tabularx}
\vspace{2pt}
\parbox{\columnwidth}{\footnotesize $n_{10}$: correct by baseline, incorrect by Ensemble; $n_{01}$: incorrect by baseline, correct by Ensemble. $\alpha_{\text{Holm}}$ denotes step-down threshold for $\alpha = 0.05$.}
\end{table}

The Ensemble significantly outperformed Swin Transformer ($\chi^2 = 31.53, p < 10^{-4}$). The comparison with DINO ViT yielded $p = 0.0528$, which did not reach the prespecified significance threshold under raw or Holm-adjusted criteria ($\alpha_{\text{Holm}} = 0.0167$). Differences between ConvNeXt, Hybrid Fusion, and the Ensemble were not statistically significant under the applied test, indicating that their primary-test performance was statistically comparable.

\subsection{Ablation Analysis: Objective Formulations and Feature Branches}
\label{subsec:ablation}

To isolate the individual contribution of the distance-aware ordinal penalty, Table~\ref{tab:loss_ablation} presents a controlled loss formulation ablation on the primary held-out test partition ($N = 1{,}656$, 828 patients) under identical training schedules, optimizers, and seeds.

\begin{table}[t]
\centering
\caption{Ablation of loss formulations on the primary held-out test partition ($N = 1{,}656$, 828 patients) using the Hybrid Fusion architecture under static inference. Extreme reversals denote Severe (KL~4) instances mispredicted as Low/No OA (KL~0--1).}
\label{tab:loss_ablation}
\scriptsize
\setlength{\tabcolsep}{3pt}
\renewcommand{\arraystretch}{1.1}
\begin{tabularx}{\columnwidth}{@{}Xccccc@{}}
\toprule
\textbf{Optimization Objective} & \textbf{Acc (\%)} & \textbf{F1 (\%)} & \textbf{QWK} & \textbf{Sev Rec (\%)} & \textbf{Reversals} \\
\midrule
Standard Cross-Entropy (CE, $\lambda=0$) & 81.28 & 80.85 & 0.678 & 94.12 & 2 \\
CE + Static Cost-Sensitive Matrix & 82.12 & 81.90 & 0.697 & 96.08 & 1 \\
CE + Chen \textit{et al.} Ordinal Loss \cite{chen2019fully} & 82.79 & 82.64 & 0.709 & 98.04 & 0 \\
Proposed ($\lambda_{\text{ord}} = 0.25$) & 83.03 & 82.91 & 0.714 & 98.04 & 0 \\
\textbf{Proposed ($\lambda_{\text{ord}} = 0.50$, Default)} & \textbf{83.44} & \textbf{83.38} & \textbf{0.721} & \textbf{100.00} & \textbf{0} \\
Proposed ($\lambda_{\text{ord}} = 1.00$) & 82.61 & 82.40 & 0.724 & 100.00 & 0 \\
\bottomrule
\end{tabularx}
\end{table}

The ablation provides evidence that distance-aware optimization improves ordinal agreement and suppresses observed non-adjacent errors: standard cross-entropy achieves 81.28\% accuracy and 0.678 QWK, but incurs 2 catastrophic non-adjacent reversals (Severe misdiagnosed as Low/No OA). The static cost-sensitive matrix baseline incorporates a fixed misclassification cost matrix $\mathbf{C} \in \mathbb{R}^{3 \times 3}$ with linear step penalties: $\mathbf{C} = \begin{bmatrix} 0 & 1 & 2 \\ 1 & 0 & 1 \\ 2 & 1 & 0 \end{bmatrix}$ (where $C_{ij} = |i - j|$), weighting prediction penalties by discrete class distance. While static cost weighting elevates QWK to 0.697 and improves Severe recall to 96.08\%, it lacks distribution-level probability mass shaping and still permits 1 extreme reversal. In contrast, the proposed distance-aware formulation ($\lambda_{\text{ord}} = 0.50$) improves accuracy by +2.16 pp to 83.44\%, elevates QWK by +0.043 to 0.721, achieves 100\% Severe recall on held-out test data, and completely eliminates extreme grade reversals. During hyperparameter tuning strictly on the validation partition ($\mathcal{P}_{\text{Val}}$, $N = 826$ radiographs across 413 patients), candidate ordinal penalty weights were evaluated prior to touching the held-out test data: $\lambda_{\text{ord}} = 0.00$ (CE only: 81.72\% val accuracy, 0.684 val QWK, 1 Severe omission into Low/No OA), $\lambda_{\text{ord}} = 0.25$ (83.17\% val accuracy, 0.718 val QWK, 0 omissions), $\lambda_{\text{ord}} = 0.50$ (83.66\% val accuracy, 0.729 val QWK, 0 omissions), and $\lambda_{\text{ord}} = 1.00$ (82.81\% val accuracy, 0.726 val QWK, 0 omissions). Maximizing validation QWK while eliminating non-adjacent omissions identified $\lambda_{\text{ord}} = 0.50$ as the optimal trade-off, which was frozen for all subsequent test evaluations. The post-hoc held-out test ablation in Table~\ref{tab:loss_ablation} indicates that this validation-derived configuration generalizes robustly to unseen test radiographs, advancing accuracy by +2.16 pp and QWK by +0.043 over standard cross-entropy while completely eliminating extreme reversals.

\begin{table}[t]
\centering
\caption{Fixed-head branch sensitivity analysis on the Hybrid architecture ($N = 1{,}656$, static inference). Feature outputs of the specified branch were zeroed in the fused latent representation $\mathbf{z}_{\text{fused}}$ prior to evaluation.}
\label{tab:branch_ablation}
\scriptsize
\setlength{\tabcolsep}{4pt}
\renewcommand{\arraystretch}{1.1}
\begin{tabularx}{\columnwidth}{@{}Xccc@{}}
\toprule
\textbf{Ablated Architectural Branch} & \textbf{Accuracy (\%)} & \textbf{QWK} & \textbf{$\Delta$ Accuracy (pp)} \\
\midrule
None (Complete Hybrid Baseline) & 83.44 & 0.721 & --- \\
ConvNeXt Branch Zeroed (Local Texture) & 76.12 & 0.615 & -7.32 \\
Swin Branch Zeroed (Global Relational) & 78.45 & 0.648 & -4.99 \\
DINO ViT Branch Zeroed (Structural Geometry) & 75.80 & 0.608 & -7.64 \\
\bottomrule
\end{tabularx}
\vspace{2pt}
\parbox{\columnwidth}{\footnotesize $\Delta$ Accuracy denotes absolute change in percentage points relative to the unablated Hybrid baseline (83.44\% accuracy, 0.721 QWK).}
\end{table}

Table~\ref{tab:branch_ablation} details the fixed-head sensitivity analysis conducted by zeroing specific backbone representations in $\mathbf{z}_{\text{fused}}$. Zeroing the self-supervised DINO ViT representation produced the largest performance drop ($-7.64$ pp accuracy, $-0.113$ QWK), followed closely by ConvNeXt ablation ($-7.32$ pp accuracy). These drops suggest that self-supervised distilled geometric boundaries and localized trabecular texture contribute substantial complementary evidence within the fused feature space.

\subsection{Qualitative Visual Explainability}
\label{subsec:xai_results}

\begin{figure*}[!t]
    \centering
    \includegraphics[width=0.95\textwidth]{fig4.jpeg}
    \caption{Multi-method qualitative explainability comparison across six attribution algorithms evaluated on representative radiographs spanning all three diagnostic tiers (Row 1: Low/No OA KL~0--1; Row 2: Moderate OA KL~2--3; Row 3: Severe OA KL~4). Columns from left to right: (Col 1) Original weight-bearing radiograph, (Col 2) Vanilla Saliency map, (Col 3) Grad-CAM localization, (Col 4) Grad-CAM++ capturing multi-focal features, (Col 5) Layer-CAM preserving fine-grained spatial activations, (Col 6) Score-CAM gradient-free forward-pass weighting, and (Col 7) Axiomatic Integrated Gradients satisfying path-integral completeness. Heatmaps demonstrate consistent qualitative localization along tibiofemoral joint margins, subchondral sclerosis, and osteophytes without spurious background activations.}
    \label{fig:xai}
\end{figure*}

Figure~\ref{fig:xai} presents a qualitative visual explainability comparison across six gradient- and activation-based attribution methods evaluated on representative correctly classified radiographs across Low/No OA (Row 1), Moderate OA (Row 2), and Severe OA (Row 3). Attribution methods operate on the terminal ConvNeXt convolutional stage and input space to visualize morphological features utilized during inference. To prevent cherry-picking, candidate test cases were sampled randomly from the primary held-out test cohort across correctly predicted cases in each diagnostic tier, ensuring that visualized patterns reflect characteristic morphological hallmarks of joint degeneration rather than idiosyncratic visual artifacts.

In Vanilla Saliency (Col 2), pixel-level gradient backpropagation highlights high-frequency subchondral bone contours along the medial tibial plateau. Grad-CAM (Col 3) and Layer-CAM (Col 5) delineate the broader articular interface: in Low/No OA, activations distribute across the preserved medial and lateral joint compartments, whereas in Moderate and Severe cases, heatmaps concentrate over the compromised medial compartment where joint space loss is pronounced. Grad-CAM++ (Col 4) and Score-CAM (Col 6) demonstrate granular localization over osteophytic spurs protruding from the tibial spine. Integrated Gradients (Col 7) shows that attribution localizes over subchondral sclerosis and joint margin narrowing. Across the examined instances, attribution maps are qualitatively consistent with recognized anatomical hallmarks of degenerative joint disease.

\begin{figure*}[!b]
    \centering
    \includegraphics[width=0.78\textwidth]{fig5.jpeg}
    \caption{Model-agnostic Local Interpretable Model-agnostic Explanations (LIME) superpixel interpretability across diagnostic tiers (Row 1: Low/No OA, Row 2: Moderate OA, Row 3: Severe OA). Columns from left to right: (Col 1) Original knee radiograph, (Col 2) SLIC superpixel segmentation partitioning joint anatomy into spatial segments, (Col 3) LIME local ridge regression feature importance heatmap, and (Col 4) Explanatory region overlay, where green contours designate superpixels providing positive evidence supporting the predicted tier, while red contours denote contradictory evidence.}
    \label{fig:lime}
\end{figure*}

To complement gradient-based saliency with a perturbation-based technique, Figure~\ref{fig:lime} illustrates Local Interpretable Model-agnostic Explanations (LIME) across four procedural stages: original radiograph (Col 1), Simple Linear Iterative Clustering (SLIC) superpixel partitioning (Col 2), local ridge regression feature importance map (Col 3), and explanatory region overlays (Col 4). In Low/No OA (Row 1), positive superpixels encircle the patent inter-articular cartilage gap. In Moderate OA (Row 2), positive superpixels cluster along medial tibial spurs. In Severe OA (Row 3), positive green superpixels isolate medial joint space obliteration and dense subchondral sclerosis. The representative attribution maps were qualitatively consistent with established radiographic signs of joint degeneration across both gradient attributions (Fig.~\ref{fig:xai}) and input perturbations (Fig.~\ref{fig:lime}).

\subsection{Computational Complexity and Inference Latency Benchmarking}
\label{subsec:computational_benchmarking}

To evaluate translational viability, Table~\ref{tab:computational_complexity} summarizes the parameter volume, floating-point operations (FLOPs), and hardware execution latency across all architectural configurations on standardized $224 \times 224 \times 3$ radiographs.

\begin{table}[t]
\centering
\caption{Computational complexity, parameter volume, and inference latency profiling across individual backbones and the hybrid architecture evaluated on standardized $224 \times 224 \times 3$ radiographs.}
\label{tab:computational_complexity}
\scriptsize
\setlength{\tabcolsep}{3pt}
\renewcommand{\arraystretch}{1.1}
\begin{tabularx}{\columnwidth}{@{}Xcccc@{}}
\toprule
\textbf{Architecture} & \textbf{Parameters (M)} & \textbf{FLOPs (G)} & \textbf{Latency (ms/img)} & \textbf{Throughput (img/s)} \\
\midrule
ConvNeXt-Base & 87.6 & 15.4 & $4.6 \pm 0.3$ & 217.4 \\
Swin-Base & 86.7 & 15.4 & $5.1 \pm 0.4$ & 196.1 \\
DINO ViT-Base & 85.8 & 17.6 & $5.8 \pm 0.4$ & 172.4 \\
MLP Projection Head & 3.41 & 0.01 & $0.3 \pm 0.05$ & 3{,}333.3 \\
\midrule
\textbf{Full Hybrid (Static)} & \textbf{263.5} & \textbf{48.4} & $\mathbf{14.2 \pm 0.8}$ & \textbf{70.4} \\
\textbf{Full Hybrid (+ TTA)} & \textbf{263.5} & \textbf{96.8} & $\mathbf{26.8 \pm 1.2}$ & \textbf{37.3} \\
\bottomrule
\end{tabularx}
\vspace{2pt}
\parbox{\columnwidth}{\footnotesize Parameter counts report actual instantiated module parameters $\sum p.\text{numel()}$ with unused pretraining/classification heads removed. Latency denotes per-image inference time ($T_{\text{batch}} / B$ with batch size $B = 16$) measured on an NVIDIA GeForce RTX 4050 Laptop GPU (6~GB GDDR6, CUDA 12.1, PyTorch 2.4.0) under mixed precision (FP16) after 10 warmup iterations, averaged over 100 synchronized runs. Throughput is calculated consistently as $1000 / \text{Latency (ms/img)}$.}
\end{table}

The tri-paradigm Hybrid architecture consolidates 263.5 million total instantiated parameters: 87.6M allocated to ConvNeXt-Base (87,566,464), 86.7M to Swin-Base (86,743,224), 85.8M to frozen DINO ViT-Base (85,798,656), and 3.41M to the dual-stage MLP projection head ($2816 \to 1024 \to 512 \to 3$ with Batch Normalization, totaling 3,414,019 parameters). In Stage~1, only the 3.41M MLP head parameters (3,414,019) are optimized. During Stage~2 selective fine-tuning, earlier backbone stages and the entirety of DINO remain frozen, updating only the terminal stages of ConvNeXt (\texttt{stages.3}, 27.44M) and Swin (\texttt{layers.3}, 27.30M) alongside the head, resulting in 58.2 million active trainable parameters (58,161,731). Forward execution across the three parallel backbones incurs 48.4 GFLOPs per single-pass radiograph, evaluated under the standard convention of 2 floating-point operations per multiply-accumulate (totaling 24.2 GMACs) measured via the PyTorch \texttt{thop} profiling tool. On an NVIDIA GeForce RTX 4050 GPU under mixed precision (FP16), single-pass static inference requires $14.2 \pm 0.8\text{ ms}$ per image, corresponding to a throughput of 70.4 images per second. Under horizontal-flip TTA, total computation scales to 96.8 GFLOPs, requiring $26.8 \pm 1.2\text{ ms}$ per radiograph ($37.3\text{ images/s}$). On host CPU execution (Intel Core i7-13700H @ 2.40~GHz), single-pass latency averaged $118.5 \pm 4.6\text{ ms}$. Peak GPU memory consumption during batch inference ($B=16$) remained below 2.2~GB. These measurements indicate favorable computational efficiency for potential integration into clinical decision-support workflows.

\subsection{Diagnostic Error Case Taxonomy and Borderline Boundary Analysis}
\label{subsec:error_analysis}

Auditing the specific failure modes that persist within the 4,008-image evaluation cohort indicates that classification errors concentrate entirely along adjacent clinical stages:

\begin{enumerate}[leftmargin=*]
\item \textit{Borderline Osteophytosis versus Physiological Aging (Moderate misclassified as Low/No OA):} For the Hybrid Fusion model, 234 of 1,591 Moderate radiographs (14.71\%) were assigned to Low/No OA (compared to 257 of 1,591, or 16.15\%, for the Soft-Voting Ensemble; total Moderate misclassifications across both directions were 15.08\% for Hybrid and 16.59\% for Ensemble). Retrospective qualitative review reveals that these cases predominantly represent borderline KL~2 transitions characterized by subtle marginal osteophytic spurs along the tibial spine or femoral condylar margins without definitive joint space narrowing. In these instances, the model's predicted probability distribution exhibited elevated entropy ($p_{\text{Low/No OA}} \in [0.48, 0.55]$ vs.\ $p_{\text{Moderate}} \in [0.45, 0.52]$), reflecting decision uncertainty near the boundary rather than gross diagnostic discordance.
\item \textit{Sclerotic Overestimation (Moderate misclassified as Severe):} A minimal fraction of Moderate cases were elevated to Severe: 6 cases (0.38\%) by Hybrid Fusion and 7 cases (0.44\%) by the Ensemble. Qualitative review indicates that these radiographs featured pronounced subchondral eburnation and marked spurring while preserving a narrowed but patent joint space. The intense sclerotic pixel intensities prompted the convolutional and DINO branches to signal severe joint degradation despite partial cartilage preservation.
\item \textit{Zero Severe-to-Low/No OA Reversals (0/122 instances):} In concordance with clinical triage requirements, zero Severe cases were misclassified into the Low/No OA category across all 4,008 non-training radiographs. For all 122 Severe instances, posterior probabilities for the Low/No OA tier remained below $p_0 < 0.02$, with mean predicted Severe confidence exceeding 94.6\%. No non-adjacent Severe-to-Low/No OA errors were observed in the evaluated configurations. While strong feature representations across backbones establish high baseline disease sensitivity, the controlled loss ablation provides evidence that incorporating the distance-aware ordinal penalty reduces the occurrence of observed non-adjacent errors.
\end{enumerate}


% ==============================================================================
% 5. DISCUSSION
% ==============================================================================
\FloatBarrier
\section{Discussion}
\label{sec:discussion}

\subsection{Diagnostic Efficacy and the Zero-Reversal Criterion}
\label{subsec:clinical_efficacy}

The central translational objective of automated medical image grading is to provide trustworthy triage support without introducing hazardous clinical errors. In knee osteoarthritis, mistaking an advanced joint deformity (Severe, KL~4) for a normal knee (Low/No OA, KL~0--1) represents a critical diagnostic failure that could delay specialist orthopedic evaluation \cite{carr2012knee}. Conversely, confusion between adjacent borderline grades (e.g., distinguishing early osteophytes in KL~1 vs.\ KL~2) reflects clinical ambiguity that is routinely observed among musculoskeletal radiologists \cite{sheehy2015validity}.

Our experimental results demonstrate that the proposed framework achieved 100.00\% empirical sensitivity (122/122) on the Severe cohort under dual-pass TTA. Importantly, no non-adjacent Severe-to-Low/No OA errors were observed in the evaluated configurations on the 4,008-image evaluation cohort (0/122). While high backbone discriminative capacity contributes to overall disease recognition, evidence from the controlled loss ablation (Section~\ref{subsec:ablation}) indicates that the distance-aware ordinal objective is associated with fewer severe-to-low/no-OA boundary omissions in the controlled ablation by penalizing probability mass proportionally to clinical distance. When combined with horizontal-flip test-time augmentation, prediction variance along borderline decision thresholds is reduced, yielding a Quadratic Weighted Kappa of 0.853 on the composite cohort and 0.748 on the primary held-out test partition. As delineated in Table~\ref{tab:comparison}, while direct accuracy comparisons must account for the difference between 3-class triage and 5-class nominal grading, both the primary held-out test baseline (83.44\% static, 85.08\% TTA on $N = 1{,}656$) and the expanded multi-partition evaluation (90.92\% on $N = 4{,}008$) demonstrate competitive alignment with published KOA grading benchmarks under verified patient-level leakage control.

\begin{table*}[t]
\centering
\caption{Comparison with published deep learning methods for KOA severity grading. Direct numerical comparisons must be interpreted with caution due to differing task granularities (5-class nominal KL grading vs.\ 3-class study-defined triage grouping), splitting protocols (patient-level vs.\ image-level), and imaging views. Quadratic Weighted Kappa (QWK) provides the primary chance-corrected metric for ordinal agreement across differing scales.}
\label{tab:comparison}
\scriptsize
\setlength{\tabcolsep}{4pt}
\renewcommand{\arraystretch}{1.1}
\begin{tabular}{p{3.2cm}p{2.3cm}p{1.0cm}p{2.0cm}p{2.6cm}p{1.8cm}p{1.1cm}p{1.2cm}}
\toprule
\textbf{Study} & \textbf{Architecture} & \textbf{Classes} & \textbf{Dataset} & \textbf{Evaluation Cohort} & \textbf{Accuracy} & \textbf{QWK} & \textbf{AUC} \\
\midrule
Antony \textit{et al.} \cite{antony2016quantifying} & VGG-16 & 5 & OAI & Image-level split & 57.0\% & 0.67 & --- \\
Tiulpin \textit{et al.} \cite{tiulpin2018automatic} & Siamese ResNet & 5 & OAI & Patient-level split & 66.7\% & 0.83 & --- \\
Chen \textit{et al.} \cite{chen2019fully} & Ordinal DenseNet & 5 & OAI & Patient-level split & --- & 0.79 & --- \\
Thomas \textit{et al.} \cite{thomas2020automated} & Attention CNN & 5 & OAI & Patient-level split & 71.0\% & --- & --- \\
Swiecicki \textit{et al.} \cite{swiecicki2021deep} & Dual-View CNN & 5 & MOST & Patient-level split & 71.9\% & 0.91 & --- \\
Mohammed \textit{et al.} \cite{bany2023knee} & ResNet-101 & 3 / 5 & Mendeley & Image-level split & 89.0\% / 69.0\% & --- & --- \\
Kondal \textit{et al.} \cite{kondal2023transfer} & Two-Stage CNN & Cont. & OAI/Clinical & Patient-level split & MAE = 0.28 & --- & --- \\
\midrule
\textbf{Present study (Primary Test)} & \textbf{Hybrid (Static)} & \textbf{3} & \textbf{Mendeley ($N=1{,}656$)} & \textbf{Strict Patient-Level Held-Out} & \textbf{83.44\% [81.7, 85.2]} & \textbf{0.721} & \textbf{0.943} \\
\textbf{Present study (Primary Test)} & \textbf{Hybrid + TTA} & \textbf{3} & \textbf{Mendeley ($N=1{,}656$)} & \textbf{Strict Patient-Level Held-Out} & \textbf{85.08\% [83.4, 86.8]} & \textbf{0.748} & \textbf{0.954} \\
\textbf{Present study (Stress Cohort)} & \textbf{Hybrid + TTA} & \textbf{3} & \textbf{Mendeley ($N=4{,}008$)} & \textbf{Composite Multi-Partition} & \textbf{90.92\% [90.0, 91.8]} & \textbf{0.853} & \textbf{0.980} \\
\bottomrule
\end{tabular}
\end{table*}

\subsection{Mechanistic Interpretation of Tri-Paradigm Representation}
\label{subsec:representation_mechanisms}

A central question in medical computer vision centers on whether multi-backbone fusion provides genuine representational complementarity or simply inflates parameter overhead. Our sensitivity analysis (Table~\ref{tab:branch_ablation}) indicates that each network branch isolates distinct morphological features:

\begin{enumerate}[leftmargin=*]
\item \textit{ConvNeXt-Base (Subchondral Trabecular Texture):} Depthwise separable convolutions operating over $7 \times 7$ receptive fields are intended to capture high-frequency subchondral bone texture and marginal osteophytic contours along the articular rim. Zeroing this branch produces a 7.32 percentage-point reduction in diagnostic accuracy, consistent with sensitivity to localized trabecular and margin cues.
\item \textit{Swin Transformer (Inter-Compartmental Relational Context):} Shifted-window self-attention is designed to model broader spatial dependencies across the articular gap, comparing medial-to-lateral compartment height asymmetry and overall joint alignment. Without this contextual perspective, classification accuracy drops by 4.99 pp, indicating that relational compartment representations contribute meaningful complementary signal.
\item \textit{DINO ViT-Base (Invariant Articular Geometry):} Preserving frozen self-supervised representations provides an invariant structural prior that resists overfitting on modest medical cohorts. Perturbing DINO produces the sharpest performance drop ($-7.64$ pp accuracy, $-0.113$ QWK), consistent with self-supervised visual features contributing substantial complementary structural information within the fused feature space.
\end{enumerate}

\subsection{Translational Triage Architecture and Clinical Decision Support}
\label{subsec:clinical_workflow}

Translating computer-aided detection models into radiology environments requires alignment with hospital Picture Archiving and Communication Systems (PACS) and Electronic Health Record (EHR) workflows. The proposed tri-paradigm framework is designed as a computer-aided triage filter to support prioritization of radiograph worklists:

\begin{enumerate}[leftmargin=*]
\item \textit{Tier 0 (Low/No OA, KL~0--1):} Patients exhibiting preserved joint space and absence of definitive osteophytes can be prioritized for routine non-urgent clinician review or conservative lifestyle counseling \cite{zhang2010eular}.
\item \textit{Tier 1 (Moderate OA, KL~2--3):} Radiographs with definitive osteophytosis and intermediate joint space loss can be flagged for clinician review and consideration of appropriate conservative management \cite{zhang2010eular}.
\item \textit{Tier 2 (Severe OA, KL~4):} Advanced joint collapse cases can be prioritized on specialist reading worklists to facilitate timely orthopedic consultation \cite{carr2012knee}.
\end{enumerate}

\noindent\textbf{Decision Support Notice:} The proposed model is intended strictly as computer-aided triage and decision-support information to assist clinician review. Predicted severity tiers do not constitute independent clinical diagnoses, surgical indications, or automated treatment prescriptions.

\subsection{Sociodemographic Equity and Fairness Considerations}
\label{subsec:bias_fairness}

A recognized challenge in musculoskeletal machine learning is the potential presence of demographic disparities. Prior studies have observed that underserved patient cohorts may experience substantial discordance between reported pain severity and radiographic KL grades \cite{pierson2021algorithmic}. Concurrently, knee osteoarthritis exhibits documented epidemiological disparities across biological sexes, with post-menopausal women exhibiting higher rates of medial compartment cartilage loss and functional impairment \cite{hunter2019oa, loeser2012oa}.

Single-backbone models trained with nominal cross-entropy can inadvertently exploit non-pathological correlations, such as soft-tissue thickness variations associated with patient body habitus or beam exposure discrepancies \cite{tajbakhsh2016convolutional}. In our proposed architecture, incorporating DINO's self-supervised visual representations may reduce sensitivity to non-pathological variations by emphasizing intrinsic structural contours. Qualitative explainability audits (Figs.~\ref{fig:xai} and \ref{fig:lime}) showed attribution patterns that were consistent with relevant anatomical regions of the knee joint. However, no formal demographic subgroup fairness audit was conducted in this study due to the absence of linked sociodemographic metadata in the public imaging repository. This represents an important limitation, and future multi-center evaluations must systematically validate performance across diverse patient populations.

\subsection{Methodological Limitations}
\label{subsec:critical_limitations}

Several methodological constraints should be acknowledged:
\begin{enumerate}[leftmargin=*]
\item \textit{Cohort Source and Multi-Center Generalizability:} Although the repository comprises 9,786 radiographs across 4,130 unique subjects, all underlying images originate from the Osteoarthritis Initiative (OAI) cohort \cite{nevitt2006oai}. While standardized radiographic protocols ensure rigorous internal benchmarking, external validation across independent clinical cohorts (such as the Multicenter Osteoarthritis Study [MOST], CHECK, or international institutional PACS archives) remains the essential next milestone to establish generalizability against domain shifts in detector hardware, beam exposure calibrations, and diverse patient demographics.
\item \textit{Input Resolution Downsampling:} Standardizing radiographs to $224 \times 224$ pixels ensures compatibility with pretrained vision backbones but introduces spatial downsampling that may attenuate subtle micro-structural cues, such as minute tibial spine sharpening.
\item \textit{Three-Tier Consolidation:} Grouping KL grades into Low/No OA (KL~0--1), Moderate OA (KL~2--3), and Severe OA (KL~4) reflects an operational triage grouping, but does not provide granular discrimination between KL~2 and KL~3 for monitoring multi-year structural progression.
\end{enumerate}

\subsection{Future Directions}
\label{subsec:future_scope}

Building upon these findings, future research trajectories include:
\begin{enumerate}[leftmargin=*]
\item \textit{Musculoskeletal Domain-Specific Foundation Pretraining:} Pretraining vision backbones directly on large extremity radiograph archives (such as Stanford MURA and institutional extremity PACS collections) using masked autoencoding or self-distillation to capture radiographic physics directly.
\item \textit{Multimodal Progression Modeling:} Integrating radiographic features with clinical biomarkers---such as age, body mass index, and patient-reported outcome measures (WOMAC scores)---to forecast multi-year joint collapse risk, extending preliminary progression modeling \cite{guan2022deep}.
\item \textit{External and Federated Multi-Center Clinical Validation:} Conducting prospective external validation on independent clinical cohorts (e.g., MOST, CHECK) and deploying within federated learning consortia to rigorously evaluate domain generalization across differing imaging equipment and community healthcare centers while preserving patient privacy.
\end{enumerate}


% ==============================================================================
% 6. CONCLUSION
% ==============================================================================
\section{Conclusion}
\label{sec:conclusion}

This investigation introduced an automated, leakage-controlled hybrid deep learning framework for knee osteoarthritis severity grading, synthesizing ConvNeXt-Base, Swin Transformer, and self-supervised DINO ViT through late feature fusion. Evaluated across 8,260 core radiographs from 4,130 unique subjects under verified patient-level isolation, the architecture achieves an accuracy of 83.44\%, a Quadratic Weighted Kappa of 0.721, and a macro ROC-AUC of 0.943 on the primary held-out test partition ($N = 1{,}656$, 828 patients), advancing to 85.08\% accuracy, 0.748 QWK, and 0.954 macro AUC under horizontal-flip test-time augmentation. On an expanded 4,008-image multi-partition stress cohort, the framework attained 90.92\% accuracy, 0.853 QWK, and 0.980 macro ROC-AUC, with 100\% empirical sensitivity (122/122) on Severe OA and zero observed extreme reversals into Low/No OA. Qualitative explainability analyses showed attribution patterns consistent with relevant anatomical features of knee degeneration. These findings highlight the utility of distance-aware ordinal optimization and multi-paradigm visual fusion for computer-aided radiographic triage.


% ==============================================================================
% ETHICS AND DATA USE STATEMENT
% ==============================================================================
\section*{Ethics and Data Use Statement}
This investigation represents a secondary analysis of publicly available, de-identified radiographic data from the Osteoarthritis Initiative (OAI) repository and Mendeley Data. The original OAI study obtained institutional review board approvals and written informed consent from all participants across participating clinical sites. In accordance with institutional policies for retrospective analyses of anonymized public research data, this study did not involve direct human subject contact.


% ==============================================================================
% DATA AVAILABILITY STATEMENT
% ==============================================================================
\section*{Data Availability Statement}
The knee radiograph dataset analyzed in this study is publicly accessible via Mendeley Data (DOI: 10.17632/56rmx5bjcr.1, Version~1) \cite{chen2018dataset} and sourced from the Osteoarthritis Initiative (OAI, \url{https://oai.epi-ucsf.org/datarelease}). All code for patient-level leakage verification, model training, evaluation, and explainability is available at \url{https://github.com/SettyBhavithav/Knee-OA-Diagnostic-Pipeline-Hybrid-ViT-Ensemble}.


% ==============================================================================
% REFERENCES
% ==============================================================================
\FloatBarrier
\begin{thebibliography}{51}
\bibitem{hunter2019oa}
D.~J. Hunter and S.~Bierma-Zeinstra, ``Osteoarthritis,'' \emph{The Lancet}, vol. 393, no. 10182, pp. 1745--1759, 2019.

\bibitem{loeser2012oa}
R.~F. Loeser, S.~R. Goldring, C.~R. Scanzello, and M.~B. Goldring, ``Osteoarthritis: A disease of the joint as an organ,'' \emph{Arthritis \& Rheumatism}, vol. 64, no.~6, pp. 1697--1707, 2012.

\bibitem{cross2014global}
M.~Cross, E.~Smith, D.~Hoy, S.~Nolte, I.~Ackerman, M.~Fransen, L.~Bridgett, S.~Williams, F.~Guillemin, C.~L. Hill \emph{et~al.}, ``The global burden of hip and knee osteoarthritis: Estimates from the Global Burden of Disease 2010 study,'' \emph{Annals of the Rheumatic Diseases}, vol.~73, no.~7, pp. 1323--1330, 2014.

\bibitem{woolf2003burden}
A.~D. Woolf and B.~Pfleger, ``Burden of major musculoskeletal conditions,'' \emph{Bulletin of the World Health Organization}, vol.~81, no.~9, pp. 646--656, 2003.

\bibitem{braun2012diagnosis}
H.~J. Braun and G.~E. Gold, ``Diagnosis of osteoarthritis: Imaging,'' \emph{Bone}, vol.~51, no.~2, pp. 278--288, 2012.

\bibitem{kellgren1957radiological}
J.~H. Kellgren and J.~S. Lawrence, ``Radiological assessment of osteo-arthrosis,'' \emph{Annals of the Rheumatic Diseases}, vol.~16, no.~4, pp. 494--502, 1957.

\bibitem{wright2014interobserver}
R.~W. Wright and MARS Group, ``Osteoarthritis classification scales: Interobserver reliability and arthroscopic correlation,'' \emph{The Journal of Bone and Joint Surgery}, vol.~96, no.~14, pp. 1145--1151, 2014.

\bibitem{sheehy2015validity}
L.~Sheehy, E.~Culham, L.~McLean, J.~Niu, J.~Lynch, N.~A. Segal, J.~A. Singh, M.~Nevitt, and T.~D.~V. Cooke, ``Validity and sensitivity to change of three scales for the radiographic assessment of knee osteoarthritis using images from the Multicenter Osteoarthritis Study (MOST),'' \emph{Osteoarthritis and Cartilage}, vol.~23, no.~9, pp. 1491--1498, 2015.

\bibitem{culvenor2019oa}
A.~G. Culvenor, B.~E. {\O}iestad, H.~F. Hart, J.~J. Stefanik, A.~Guermazi, and K.~M. Crossley, ``Prevalence of knee osteoarthritis features on magnetic resonance imaging in asymptomatic uninjured adults: A systematic review and meta-analysis,'' \emph{British Journal of Sports Medicine}, vol.~53, no.~20, pp. 1268--1278, 2019.

\bibitem{tiulpin2018automatic}
A.~Tiulpin, J.~Thevenot, E.~Rahtu, P.~Lehenkari, and S.~Saarakkala, ``Automatic knee osteoarthritis diagnosis from plain radiographs: A deep learning-based approach,'' \emph{Scientific Reports}, vol.~8, no.~1, p. 1727, 2018.

\bibitem{thomas2020automated}
K.~A. Thomas, L.~Kidzi{\'n}ski, E.~Halilaj, S.~L. Fleming, G.~R. Venkataraman, E.~H.~G. Oei, G.~E. Gold, and S.~L. Delp, ``Automated classification of radiographic knee osteoarthritis severity using deep neural networks,'' \emph{Radiology: Artificial Intelligence}, vol.~2, no.~2, p. e190065, 2020.

\bibitem{chen2019fully}
P.~Chen, L.~Gao, X.~Shi, K.~Allen, and L.~Yang, ``Fully automatic knee osteoarthritis severity grading using deep neural networks with a novel ordinal loss,'' \emph{Computerized Medical Imaging and Graphics}, vol.~75, pp. 84--92, 2019.

\bibitem{pierson2021algorithmic}
E.~Pierson, D.~M. Cutler, J.~Leskovec, S.~Mullainathan, and Z.~Obermeyer, ``An algorithmic approach to reducing unexplained pain disparities in underserved populations,'' \emph{Nature Medicine}, vol.~27, no.~1, pp. 136--140, 2021.

\bibitem{antony2016quantifying}
J.~Antony, K.~McGuinness, N.~E. O'Connor, and K.~Moran, ``Quantifying radiographic knee osteoarthritis severity using deep convolutional neural networks,'' in \emph{Proceedings of the 23rd International Conference on Pattern Recognition (ICPR)}. IEEE, 2016, pp. 1195--1200.

\bibitem{schiphof2008differences}
D.~Schiphof, M.~Boers, and S.~M.~A. Bierma-Zeinstra, ``Differences in descriptions of Kellgren and Lawrence grades of knee osteoarthritis,'' \emph{Annals of the Rheumatic Diseases}, vol.~67, no.~7, pp. 1034--1036, 2008.

\bibitem{zhang2010eular}
W.~Zhang, M.~Doherty, G.~Peat, S.~M.~A. Bierma-Zeinstra, N.~K. Arden, B.~Bresnihan, G.~Herrero-Beaumont, S.~Kirschner, B.~F. Leeb, L.~S. Lohmander \emph{et~al.}, ``EULAR evidence-based recommendations for the diagnosis of knee osteoarthritis,'' \emph{Annals of the Rheumatic Diseases}, vol.~69, no.~3, pp. 483--489, 2010.

\bibitem{emrani2008joint}
P.~S. Emrani, J.~N. Katz, C.~L. Kessler, W.~M. Reichmann, E.~A. Wright, T.~E. McAlindon, and E.~Losina, ``Joint space narrowing and Kellgren--Lawrence progression in knee osteoarthritis: An analytic literature synthesis,'' \emph{Osteoarthritis and Cartilage}, vol.~16, no.~8, pp. 873--882, 2008.

\bibitem{shamir2009early}
L.~Shamir, S.~M. Ling, W.~W. Scott, A.~Bos, N.~Orlov, T.~J. Macura, D.~M. Eckley, L.~Ferrucci, and I.~G. Goldberg, ``Knee X-ray image analysis method for automated detection of osteoarthritis,'' \emph{IEEE Transactions on Biomedical Engineering}, vol.~56, no.~2, pp. 407--415, 2009.

\bibitem{oka2008fully}
H.~Oka, S.~Muraki, T.~Akune, A.~Mabuchi, T.~Suzuki, H.~Yoshida, S.~Yamamoto, K.~Nakamura, N.~Yoshimura, and H.~Kawaguchi, ``Fully automatic quantification of knee osteoarthritis severity on plain radiographs,'' \emph{Osteoarthritis and Cartilage}, vol.~16, no.~11, pp. 1300--1306, 2008.

\bibitem{thomson2015automated}
J.~Thomson, T.~O'Neill, D.~Felson, and T.~Cootes, ``Automated shape and texture analysis for detection of osteoarthritis from radiographs of the knee,'' in \emph{Medical Image Computing and Computer-Assisted Intervention (MICCAI)}, 2015.

\bibitem{he2016deep}
K.~He, X.~Zhang, S.~Ren, and J.~Sun, ``Deep residual learning for image recognition,'' in \emph{Proceedings of the IEEE Conference on Computer Vision and Pattern Recognition (CVPR)}, 2016, pp. 770--778.

\bibitem{huang2017densely}
G.~Huang, Z.~Liu, L.~van~der Maaten, and K.~Q. Weinberger, ``Densely connected convolutional networks,'' in \emph{Proceedings of the IEEE Conference on Computer Vision and Pattern Recognition (CVPR)}, 2017, pp. 4700--4708.

\bibitem{tan2019efficientnet}
M.~Tan and Q.~V. Le, ``EfficientNet: Rethinking model scaling for convolutional neural networks,'' in \emph{International Conference on Machine Learning (ICML)}. PMLR, 2019, pp. 6105--6114.

\bibitem{liu2021swin}
Z.~Liu, Y.~Lin, Y.~Cao, H.~Hu, Y.~Wei, Z.~Zhang, S.~Lin, and B.~Guo, ``Swin Transformer: Hierarchical vision transformer using shifted windows,'' in \emph{Proceedings of the IEEE/CVF International Conference on Computer Vision (ICCV)}, 2021, pp. 10012--10022.

\bibitem{liu2022convnext}
Z.~Liu, H.~Mao, C.-Y. Wu, C.~Feichtenhofer, T.~Darrell, and S.~Xie, ``A ConvNet for the 2020s,'' in \emph{Proceedings of the IEEE/CVF Conference on Computer Vision and Pattern Recognition (CVPR)}, 2022, pp. 11976--11986.

\bibitem{caron2021dino}
M.~Caron, H.~Touvron, I.~Misra, H.~J{\'e}gou, J.~Mairal, P.~Bojanowski, and A.~Joulin, ``Emerging properties in self-supervised vision transformers,'' in \emph{Proceedings of the IEEE/CVF International Conference on Computer Vision (ICCV)}, 2021, pp. 9650--9660.

\bibitem{raghu2019transfusion}
M.~Raghu, C.~Zhang, J.~Kleinberg, and S.~Bengio, ``Transfusion: Understanding transfer learning for medical imaging,'' \emph{Advances in Neural Information Processing Systems (NeurIPS)}, vol.~32, 2019.

\bibitem{tajbakhsh2016convolutional}
N.~Tajbakhsh, J.~Y. Shin, S.~R. Gurudu, R.~T. Hurst, C.~B. Kendall, M.~B. Gotway, and J.~Liang, ``Convolutional neural networks for medical image analysis: Full training or fine tuning?'' \emph{IEEE Transactions on Medical Imaging}, vol.~35, no.~5, pp. 1299--1312, 2016.

\bibitem{antony2017automatic}
J.~Antony, K.~McGuinness, K.~Moran, and N.~E. O'Connor, ``Automatic detection of knee joints and quantification of knee osteoarthritis severity using convolutional neural networks,'' in \emph{Machine Learning and Data Mining in Pattern Recognition (MLDM)}. Springer, 2017, pp. 376--390.

\bibitem{swiecicki2021deep}
A.~Swiecicki, N.~Li, J.~O'Donnell, N.~Said, J.~Yang, R.~C. Mather, W.~A. Jiranek, and M.~A. Mazurowski, ``Deep learning-based algorithm for assessment of knee osteoarthritis severity in radiographs matches performance of radiologists,'' \emph{Computers in Biology and Medicine}, vol. 133, p. 104334, 2021.

\bibitem{guan2022deep}
B.~Guan, F.~Liu, A.~Haj-Mizaian, S.~Demehri, and T.~Neogi, ``Deep learning approach to predict radiographic knee osteoarthritis progression,'' \emph{Osteoarthritis and Cartilage}, vol.~27, pp. S395--S396, 2019.

\bibitem{bany2023knee}
A.~S. Mohammed, A.~A. Hasanaath, G.~Latif, and A.~Bashar, ``Knee Osteoarthritis Detection and Severity Classification Using Residual Neural Networks on Preprocessed X-ray Images,'' \emph{Diagnostics}, vol.~13, no.~8, p. 1380, 2023.

\bibitem{kondal2023transfer}
S.~Kondal, V.~Kulkarni, A.~Gaikwad, A.~Kharat, and A.~Pant, ``Automatic Grading of Knee Osteoarthritis on the Kellgren-Lawrence Scale from Radiographs Using Convolutional Neural Networks,'' in \emph{Advances in Deep Learning, Artificial Intelligence and Robotics}, Lecture Notes in Networks and Systems, vol. 249, pp. 163--173, Springer, 2022.

\bibitem{bayramoglu2021adaptive}
N.~Bayramoglu, A.~Tiulpin, J.~Hirvasniemi, M.~Niemelainen, and S.~Saarakkala, ``Adaptive segmentation of knee radiographs for selecting the optimal pretrained model for severity assessment of osteoarthritis,'' \emph{Computerized Medical Imaging and Graphics}, vol.~88, p. 101831, 2021.

\bibitem{chen2018dataset}
P.~Chen, ``Knee Osteoarthritis Severity Grading Dataset,'' \emph{Mendeley Data}, vol.~V1, 2018, doi: 10.17632/56rmx5bjcr.1.

\bibitem{carr2012knee}
A.~J. Carr, O.~Robertsson, S.~Graves, A.~J. Price, N.~K. Arden, A.~Judge, and D.~J. Beard, ``Knee replacement,'' \emph{The Lancet}, vol. 379, no. 9823, pp. 1331--1340, 2012.

\bibitem{zhang2018mixup}
H.~Zhang, M.~Cisse, Y.~N. Dauphin, and D.~Lopez-Paz, ``mixup: Beyond empirical risk minimization,'' in \emph{International Conference on Learning Representations (ICLR)}, 2018.

\bibitem{thulasidasan2019mixup}
S.~Thulasidasan, G.~Chennupati, J.~A. Bilmes, T.~Bhattacharya, and S.~Michalak, ``On mixup training: Improved calibration and predictive uncertainty for deep neural networks,'' \emph{Advances in Neural Information Processing Systems (NeurIPS)}, vol.~32, 2019.

\bibitem{loshchilov2019adamw}
I.~Loshchilov and F.~Hutter, ``Decoupled weight decay regularization,'' in \emph{International Conference on Learning Representations (ICLR)}, 2019.

\bibitem{loshchilov2017sgdr}
I.~Loshchilov and F.~Hutter, ``SGDR: Stochastic gradient descent with warm restarts,'' in \emph{International Conference on Learning Representations (ICLR)}, 2017.

\bibitem{landis1977measurement}
J.~R. Landis and G.~G. Koch, ``The measurement of observer agreement for categorical data,'' \emph{Biometrics}, vol.~33, no.~1, pp. 159--174, 1977.

\bibitem{reyes2020interpretability}
M.~Reyes, R.~Meier, S.~Pereira, C.~A. Silva, F.-M. Dahlweid, H.~von Tengg-Kobligk, R.~M. Summers, and R.~Wiest, ``On the interpretability of artificial intelligence in radiology: Challenges and opportunities,'' \emph{Radiology: Artificial Intelligence}, vol.~2, no.~3, p. e190043, 2020.

\bibitem{selvaraju2017gradcam}
R.~R. Selvaraju, M.~Cogswell, A.~Das, R.~Vedantam, D.~Parikh, and D.~Batra, ``Grad-CAM: Visual explanations from deep networks via gradient-based localization,'' in \emph{Proceedings of the IEEE International Conference on Computer Vision (ICCV)}, 2017, pp. 618--626.

\bibitem{chattopadhay2018gradcampp}
A.~Chattopadhay, A.~Sarkar, P.~Howlader, and V.~N. Balasubramanian, ``Grad-CAM++: Generalized gradient-based visual explanations for deep convolutional networks,'' in \emph{IEEE Winter Conference on Applications of Computer Vision (WACV)}, 2018, pp. 839--847.

\bibitem{jiang2021layercam}
P.-T. Jiang, C.-B. Zhang, Q.~Hou, M.-M. Cheng, and Y.~Wei, ``LayerCAM: Exploring hierarchical class activation maps for localization,'' \emph{IEEE Transactions on Image Processing}, vol.~30, pp. 5875--5888, 2021.

\bibitem{wang2020scorecam}
H.~Wang, Z.~Wang, M.~Du, F.~Yang, Z.~Zhang, S.~Ding, P.~Mardziel, and X.~Hu, ``Score-CAM: Score-weighted visual explanations for convolutional neural networks,'' in \emph{Proceedings of the IEEE/CVF Conference on Computer Vision and Pattern Recognition Workshops (CVPRW)}, 2020, pp. 24--25.

\bibitem{sundararajan2017axiomatic}
M.~Sundararajan, A.~Taly, and Q.~Yan, ``Axiomatic attribution for deep networks,'' in \emph{International Conference on Machine Learning (ICML)}. PMLR, 2017, pp. 3319--3328.

\bibitem{ribeiro2016lime}
M.~T. Ribeiro, S.~Singh, and C.~Guestrin, ``Why Should I Trust You?: Explaining the predictions of any classifier,'' in \emph{Proceedings of the 22nd ACM SIGKDD International Conference on Knowledge Discovery and Data Mining}. ACM, 2016, pp. 1135--1144.

\bibitem{johnson2019survey}
J.~M. Johnson and T.~M. Khoshgoftaar, ``Survey on deep learning with class imbalance,'' \emph{Journal of Big Data}, vol.~6, no.~1, p.~27, 2019.

\bibitem{shanmugam2021tta}
D.~Shanmugam, D.~Blalock, G.~Balakrishnan, and J.~Guttag, ``Better aggregation in test-time augmentation,'' in \emph{Proceedings of the IEEE/CVF International Conference on Computer Vision (ICCV)}, 2021, pp. 1214--1223.

\bibitem{nevitt2006oai}
M.~C. Nevitt, D.~T. Felson, and G.~Lester, ``The Osteoarthritis Initiative: Protocol for the cohort study,'' University of California San Francisco, Technical Report, 2006. [Online]. Available: \url{https://oai.epi-ucsf.org/}
\end{thebibliography}

\end{document}
